\chapter{开放世界空间 AI}
\label{ch:openworld}

% ════════════════════════════════════════════════════════════════
%  新部「学习时代的 SLAM 与空间 AI」（内容上的）收官章。
%  吸收 SLAM Handbook ch17《Towards Open-World Spatial AI》(Paull, Morin,
%    Maggio, Büchner, Cadena, Valada, Carlone)：
%    §17.1-2 基础模型(CLIP/SAM/ViT/LLM/VLM) → 压成最小必要背景；
%    §17.3 开放词汇建图（硬核、自包含主体）：range-feature 观测、稠密/物体/
%      隐式三类表示、零样本查询；Clio 信息瓶颈（式17.8/17.9）作数学重点，
%      接本书既有信息论语言（eq:mutual-info / eq:gauss-mutual-info / eq:gauss-kl，
%      ch:slam_est 的主动感知）；
%    §17.4 VLA / 「是否还需要显式地图」→ 展望/结语素材，首尾呼应 preface 的
%      具身/VLA 承诺与 intro sec:intro-need「是否需要 SLAM」。
%  最高准则：全书「经典→学习→开放世界」弧线的落点，水到渠成的收束。
%  独立性红线：§17.4 是 ventriloquize 重灾区——剥立场、只留事实 + 本书口吻综述，
%    严禁「某人认为/坦言/we argue」。
%  右扰动为主线；语义特征记 \boldsymbol\psi（严格避让 dense 章视线 \mathbf{f}、
%    learning 章光流 \mathbf{f}^flow）；地图元素 m_k=(\mathbf{p}_k,\boldsymbol\psi_k)；
%    互信息沿用 ch:slam_est 的 I(·,·)。
%  兄弟章：学习赋能（ch:learning）、动态可变形（ch:dynamic）、度量-语义
%    （ch:semantic，闭集语义，过门处直接接住其闭集 punt）、结语（ch:epilogue）。
%  教学层遵循 v5.0 算法工程教学规范，章结构对齐 lidar_slam.tex 范本。
% ════════════════════════════════════════════════════════════════

\begin{practice}[前置自测]
开始之前，先试答下列问题。它们勾勒了本章——也是\emph{学习时代}这一部走到尽头要回答的核心疑问。若已能流畅作答，可只速读本章的框架图与速查表；若卡住，正是本章要为你打通的。
\begin{enumerate}
  \item 上一章（\cref{ch:semantic}）的语义地图，把每个地图元素贴上“桌子/门/墙”之类的\emph{类别标签}。可如果机器人收到的指令是“找一把\emph{适合一边喝咖啡一边看书}的椅子”，或“拿那幅\emph{画着罗纳尔多}的画”——预先定好的类别词典还够用吗？缺了什么？
  \item 假如不再往地图里存\emph{标签}，而是存一段\emph{高维向量}，并要求“任意一句话也能编码成同一空间里的向量、谁离得近就是谁”。这种地图凭什么能回答\emph{建图时根本没预想过}的查询？它和闭集语义地图的本质差别在哪一步？
  \item “开放词汇前端”（如能接受任意文本提示的检测器）摆在 SLAM 前端，得到的地图就是“开放世界”的吗？还是说，只要\emph{建图那一刻}就把要找的概念定死了，地图骨子里仍是闭集？
  \item 同一个房间，要“搬钢琴”时它是\emph{一个}物体，要“弹钢琴”时它是\emph{九十多个}琴键，要“调音”时它又是\emph{上百根}琴弦——到底该把场景切成几个物体，谁说了算？这个看似哲学的问题，能不能写成一个\emph{优化问题}？
  \item 既然大模型能直接“看”几十帧图像回答空间问题，那还要不要费劲维护一张显式地图？这个问题，和本书开篇 \cref{sec:intro-need} 追问的“到底要不要 SLAM”，是同一个问题吗？
\end{enumerate}
\end{practice}

\section*{本章目标}
学完本章，你应当能够：
\begin{enumerate}
  \item \textbf{说清}从\emph{闭集语义}（\cref{ch:semantic}）到\emph{开集/开放词汇}的这一步跨越：闭集地图把“是什么”压成有限词典里的一个标签，开放世界地图改存基础模型的\emph{连续特征}，从而支持\emph{建图时未知}的查询；并能辨析“开放词汇前端”与“开放世界地图”两种范式的根本区别；
  \item \textbf{复述}读懂本章所需的三块基础模型背景——特征/嵌入与余弦相似度、多模态对齐（CLIP 的共享嵌入空间）、类无关分割（SAM）——并解释为什么 CLIP 的“图文共享空间”是后文一切开放词汇查询的支点；
  \item \textbf{推导}开放词汇建图的统一框架：把地图元素写成“几何 $+$ 特征” $m_k=(\mathbf{p}_k,\boldsymbol\psi_k)$、把观测写成与 \cref{ch:semantic} 的 range-category 平行的 \emph{range-feature 观测}、把跨帧特征融合写成置信加权的递归更新（\cref{eq:ow-fusion}）、把零样本查询写成特征空间里的余弦近邻（\cref{eq:ow-query}）；
  \item \textbf{比较}稠密、物体级 / 场景图、隐式三类开放词汇表示的取舍（粒度、内存、可查询性），并说清它们\emph{共享}同一套“特征嵌入 $+$ 余弦查询”内核、只在“特征存在哪一级几何上”分岔；
  \item \textbf{掌握}本章的数学重点——\emph{任务驱动的地图压缩}：用经典\textbf{信息瓶颈}把“场景该切成几个物体”形式化为 \cref{eq:ow-ib}，理解它如何把本书的\emph{互信息}（\cref{eq:mutual-info}）与“鲁棒感知时代”的\emph{任务驱动感知}（\cref{tab:intro-eras}）缝在一起，并看懂用凝聚式信息瓶颈（\cref{eq:ow-aib}）增量求解的思路；
  \item \textbf{综述}“是否还需要显式地图”这一开放问题的\emph{事实面}：长上下文 VLM 直喂图像帧、机器人基础模型与 VLA（视觉–语言–动作）的数据与模型谱系，并把它\emph{接回}本书开篇 \cref{sec:intro-need} 的“是否需要 SLAM”，给出本书自己的判断。
\end{enumerate}

\subsection*{本章知识导航}
本章是全书“经典 $\to$ 学习 $\to$ 开放世界”弧线（\cref{sec:intro-frontier}、前言）在内容上的\textbf{落点}。前面整部 SLAM 把地图做成\emph{几何}（\cref{ch:dense}）；上一章（\cref{ch:semantic}）给几何叠上\emph{闭集语义}——但语义词典在建图时就被钉死了。本章问的是最后一步：\emph{如何让地图回答任何概念，哪怕它在建图时从未被预想？}答案是把闭集标签换成基础模型的\emph{开放词汇特征}。我们沿“为什么要开集 $\to$ 用什么模型 $\to$ 怎么把特征织进地图 $\to$ 任务驱动该切多细 $\to$ 还要不要地图”五步推进，最后一步与本书开篇首尾相扣。结构如下：

\begin{center}
\begin{forest} navtreeLR
[{闭集语义\\$m_k=(\mathbf{p}_k,\sigma_k)$}
  [{开放词汇\\$m_k=(\mathbf{p}_k,\boldsymbol\psi_k)$}
    [{基础模型\\CLIP / SAM}]
    [{稠密\\range-feature}
      [{物体 / 场景图\\一物一向量}]
      [{隐式\\LERF / LangSplat}]
      [{任务驱动\\信息瓶颈 Clio}
        [{地图 API\\接 LLM}]
        [{VLA / 还要地图吗\\呼应 \cref{sec:intro-need}}]
      ]
    ]
  ]
]
\end{forest}
\end{center}

\noindent\textbf{推荐路径}：\cref{sec:openworld-closed-open}（从闭集到开集的跨越，并立“空间 AI”工作定义）是全章的“锚”，任何读者都应先读；\cref{sec:openworld-fm}（三块基础模型背景）门槛最低，只取读懂后文所需的最小必要内容；\cref{sec:openworld-mapping}（开放词汇建图）是本章技术主体，含 range-feature 观测、特征融合与零样本查询的完整推导，以及稠密/物体/隐式三类表示的对照；\cref{sec:openworld-task}（任务驱动表示与信息瓶颈）是全章\emph{理论核心}，把“场景该切多细”写成信息瓶颈并接上本书信息论语言；\cref{sec:openworld-future}（是否还需要地图 / VLA）是\emph{展望与收束}，与开篇呼应。带（选读）标记的隐式表示细节、地图 API、VLA 模型谱系面向研究，可选深潜。

\subsection*{前置知识桥接}
本章把这一部前几章的“学习层”视角推到尽头，请先激活以下前置：
\begin{itemize}
  \item \textbf{闭集度量-语义}（\cref{ch:semantic}）：上一章把语义当成有限词典里的\emph{类别}，写出 range-category 观测与多类贝叶斯融合。本章\emph{接着它的边界}讲：当词典再大也不够用时，怎么办。读懂上一章的“类别标签 $\sigma$”，才看得清本章“特征向量 $\boldsymbol\psi$”替换了什么。
  \item \textbf{二值/多类占据融合}（\cref{sec:dense-octomap} 的 \cref{eq:dense-octo-logodds}；\cref{ch:semantic} 的多类 log-odds）：本章的 range-feature 观测与特征融合（\cref{eq:ow-fusion}）是这一脉络在\emph{连续特征}上的延伸——从“逐体素一个占据概率”到“逐元素一个类别分布”再到“逐元素一个特征向量”，是同一条\emph{逐元素证据累积}主线。
  \item \textbf{学习层接母方程}（\cref{ch:learning} 的 \cref{tab:learn-spectrum}）：本部开篇 \cref{ch:learning} 立下的统一立场——“学习不推翻 MAP 母方程（\cref{eq:intro-nls}），只换其零件”。本章把这一立场推到\emph{状态表示}与\emph{语义层}：地图元素的几何 $\mathbf{p}_k$ 仍由前面各章的 SLAM 后端解出，本章只往上嵌一层\emph{特征}。几何与语义\emph{解耦}，正是开放词汇建图能站在经典 SLAM 肩上的前提。
  \item \textbf{互信息、熵与 KL 散度}（\cref{sec:est-mahalanobis} 的 \cref{eq:mutual-info,eq:gauss-mutual-info,eq:gauss-kl}）：\textbf{本章数学核心的命脉前置}。本书早已在状态估计章用本书记号定义了互信息 $I(\mathbf{x},\mathbf{y})$、并用它刻画“主动感知 / 下一最佳观测”。本章的\emph{信息瓶颈}（\cref{eq:ow-ib}）正是建在同一个 $I(\cdot,\cdot)$ 之上——读过那里，本章的“任务驱动地图压缩”就不是空降的新理论，而是同一信息论语言的又一次落地。
  \item \textbf{是否需要 SLAM}（\cref{sec:intro-need}）：开篇引 Cadena 等追问过“自主机器人真的需要 SLAM 吗”。本章末（\cref{sec:openworld-future}）在大模型时代\emph{重新追问}同一问题——这是全书首尾相扣的结构落点，读本章末前请回看开篇那一节。
  \item \textbf{相机模型与反投影}（\cref{ch:camera} 的 $\boldsymbol\pi,\boldsymbol\pi^{-1}$）：本章把像素特征抬进三维仍用针孔反投影，与 \cref{sec:dense-rgbd-pc} 由 RGB-D 生成有序点云的流程一致。
\end{itemize}

\subsection*{如果跳过本章会怎样}
\begin{itemize}
  \item 你会看到 ConceptFusion、ConceptGraphs、HOV-SG、Clio、LERF/LangSplat 这些“能用自然语言查地图”的系统，却以为它们是与本书前面 SLAM \emph{无关}的“另一套大模型应用”，从而读不懂它们其实仍\emph{站在经典 SLAM 的位姿之上}、只在几何地图上多嵌了一层特征；
  \item 你会错过全书“经典 $\to$ 学习 $\to$ 具身/开放世界”弧线的\emph{收束}：前言承诺的“通向现代具身智能的桥梁”、开篇 \cref{sec:intro-need} 埋下的“是否需要 SLAM”，都在本章末与 VLA、长上下文 VLM 的事实合流，给出本书自己的回答。跳过本章，这条弧线就缺了落点。
\end{itemize}

\begin{note}
\textbf{本章记号与约定.}\;
本章在前几部记号上沿用并扩展“学习/语义层”符号，并\emph{严格避让}全书已占用者，登记如下。

\textbf{特征 / 嵌入}统一记 $\boldsymbol\psi\in\mathbb{R}^{D}$（embedding/feature/code 的统称；\textbf{特意不用 $\mathbf{f}$}——本书 \cref{ch:dense} 已用 $\mathbf{f}$ 表“单位视线方向”、\cref{ch:learning} 已用 $\mathbf{f}^{\mathrm{flow}}$ 表光流场，故本章语义特征一律记 $\boldsymbol\psi$ 以避歧义）。\textbf{编码器}记 $\mathcal{E}_{X}$（模态 $X$ 的编码器，把输入 $x$ 映成特征 $\boldsymbol\psi_x=\mathcal{E}_X(x)$）；常用下标 $\mathcal{E}_{\mathrm{img}}$（图像）、$\mathcal{E}_{\mathrm{txt}}$（文本）。两特征的\emph{余弦相似度}记 $\langle\boldsymbol\psi_a,\boldsymbol\psi_b\rangle$（默认 $\boldsymbol\psi$ 已单位化，故内积即余弦）。

\textbf{地图元素}记 $m_k=(\mathbf{p}_k,\boldsymbol\psi_k)$：几何分量 $\mathbf{p}_k$（体素位置 / 三维点 / 网格面，\emph{由上游 SLAM 给出}）$+$ 特征分量 $\boldsymbol\psi_k$；地图 $\mathcal{M}=\{m_k\}$。\textbf{类无关分割}记 $\mathrm{Seg}(I)=\{Z_i\}$（一幅图 $I$ 上的若干掩膜，$Z_i\in\{0,1\}^{H\times W}$）。\textbf{文本/查询提示}记 $T$、其特征 $\boldsymbol\psi_T=\mathcal{E}_{\mathrm{txt}}(T)$；候选类别集记 $\mathsf{C}$。

\textbf{物体}记 $o_j=(\mathsf{P}_{o_j},\boldsymbol\psi_{o_j},n_{o_j})$（点云 $+$ 特征 $+$ 观测计数）；物体集 $\mathcal{V}_O=\{o_j\}$；几何/语义相似度记 $\phi_{\mathrm{geo}},\phi_{\mathrm{sim}}\in[0,1]$。\textbf{分层场景图}记 $\mathcal{G}=(\mathcal{V},\mathcal{E})$（与因子图、\cref{ch:semantic} 的场景图同记号）。

\textbf{信息瓶颈}：原始基元集记 $\mathcal{X}$、压缩表示记 $\tilde{\mathcal{X}}$、任务集记 $\mathcal{Y}$；软分配 $p(\tilde{x}\mid x)$；权衡系数记 $\beta$。\textbf{互信息} $I(\cdot,\cdot)$、\textbf{Jensen--Shannon 散度} $D_{\mathrm{JS}}$ 沿用 \cref{sec:est-mahalanobis} 的信息论记号（\cref{eq:mutual-info,eq:gauss-kl}）。

\textbf{相机位姿}沿用 $\mathbf{T}\in\SE(3)$ 与本书帧下标 $\mathbf{T}_{ab}$（把 $b$ 系坐标变到 $a$ 系；Handbook 原文上下标式一律改写为本书下标式，见 \nameref{ch:notation}）；扰动以\emph{右扰动}为主线。本章几乎不显式动到位姿求导（开放词汇建图在\emph{给定位姿}下进行），但凡涉及处一律按右扰动约定，与全书一致。
\end{note}

% ════════════════════════════════════════════════════════════════
\section{从闭集到开集：最后一步跨越}
\label{sec:openworld-closed-open}

\paragraph{动机：当词典再大也不够用。}
回看这一部走过的路。\cref{ch:dense} 把地图做成几何——体素的占据、TSDF 的距离、面元的表面。\cref{ch:semantic} 往几何之上叠了一层\emph{语义}：每个地图元素不仅有“在哪”，还有“是什么”——一个取自有限词典的\emph{类别标签} $\sigma$（桌子、门、墙）。这层语义是真实的进步：它让机器人能听懂“去厨房”而非只会“走到坐标 $(3.2,1.7)$”。但它有一个在建图那一刻就被钉死的前提——\textbf{闭集假设}（closed-set assumption）：地图能表达的概念，被前端感知模型（分类器、检测器）训练时的\emph{固定词典}框死了\cite{paull2026openworld}。

这个前提的代价，用一个工具箱就能看穿\cite{paull2026openworld}。箱子里有锤子、螺丝刀、扳手。一个普通分类器也许只会贴一个标签“工具”——可我们想要的是“螺丝刀”“锤子”。再细一点：这把锤子是\emph{红色}的、那把螺丝刀是\emph{蓝色}的；这件是\emph{金属}的、那件是\emph{橡胶}的；对机器人尤其要紧的是它们的\emph{用途与可供性}（affordance）——“用来切割”“用来测量”。还有墙上那幅画，理想的地图既要知道它是“一幅画”（媒介），又要知道画的是“罗纳尔多”（内容，甚至是个特定文化指称）。想用一个预定义的类别集同时满足这些，是一件苦差事：词典哪怕含上百类，对常见物体也会在某个层面失灵，对罕见的“长尾”物体更是一开始就不在词典里\cite{paull2026openworld}。还有些连续属性——一个物体\emph{典型会发出什么声音}——本就难用离散标签表达。

\begin{insight}[一句话：开集要换掉的，是“标签”这一步]
闭集语义把“是什么”压成有限词典里的\emph{一个标签} $\sigma$；这一压缩在建图时就丢掉了词典之外、以及标签粒度之下的一切信息。\textbf{开放世界建图的全部要义，是把这一步“贴标签”换成“存一段高维连续特征 $\boldsymbol\psi$”}——特征由在\emph{互联网级}数据上训练的基础模型给出，因而能表征一个\emph{开放}的概念集；而“它是不是某概念”的判断，\emph{推迟到查询时}用特征空间里的相似度来回答。一句话：闭集在\emph{建图}时定死了能问什么，开集把这件事推迟到\emph{用图}时。
\end{insight}

\paragraph{两种“开放”，只有一种是真开放。}
这里有一个极易混淆、却最关键的概念区分\cite{paull2026openworld}。要把“开集”用进 SLAM，至少有两条路：

\begin{itemize}
  \item \textbf{开放词汇\emph{前端}}：把前端感知模型换成能接受任意文本提示的“开放词汇”版本——例如 YOLO-World（YOLO 的开放词汇变体）或 OV-DETR（DETR 的开放词汇变体）能根据任意给定的文本去检测物体。\emph{但}：这种文本条件化必须发生在\emph{查询感知模型的那一刻}，在 SLAM 里就是\emph{建图的那一刻}。一旦那时定下了“要找哪些概念”，得到的地图骨子里仍是\textbf{闭集}——能表达什么，在建图时就被这次提示锁死了。
  \item \textbf{开放词汇\emph{地图}}：把开放词汇\emph{特征}嵌进地图本身，使得\emph{用图}时才提出查询、且查询可以是建图时从未预想的任何概念。\textbf{只有这一条，才是真正的开放世界。}
\end{itemize}

\noindent 这条区分是本章的判别准绳：判断任何“开放词汇 SLAM”是否名副其实，只需问——\emph{它把概念集定死在建图时，还是把判断推迟到查询时？}本章 \cref{sec:openworld-mapping} 起，一律聚焦后者。

\begin{note}[“空间 AI”：本书的工作定义]
本章标题里的“空间 AI”（Spatial AI）是个被广泛使用、却少有正式定义的术语。本书不强立一个权威定义，但给出一个贯穿本章的\emph{工作内涵}：所谓空间 AI，是指地图\emph{超越纯几何}、进一步\textbf{编码语义、并支持更抽象的高层推理}的能力\cite{paull2026openworld}。\cref{ch:dense} 的几何地图是其“地”，\cref{ch:semantic} 的闭集语义是其“阶”，本章的开放词汇与任务驱动表示是其“顶”——目标是让机器人能用\emph{直觉层面}的指令（“拿那个红色的、能喝的东西”）而非\emph{低层几何}指令（“移动到 $(x,y,z)$”）来工作。这一诉求并非凭空而来：它正是 \cref{tab:intro-eras} 所述“鲁棒感知时代”四诉求里的\emph{高层理解}与\emph{任务驱动感知}——本章是那张时代表上最后两格诉求的兑现。
\end{note}

\paragraph{一个贯穿全章的统一框架。}
为不把开放世界讲成一堆系统名的并列清单，我们沿用 \cref{ch:semantic} 已经立好的“地图元素”框架，只改其语义分量。把地图写成一组元素的集合（不失一般性）：
\begin{equation}
  \mathcal{M}\triangleq\{m_k\},\qquad m_k\triangleq(\mathbf{p}_k,\ \boldsymbol\psi_k),
  \label{eq:ow-mapelem}
\end{equation}
其中 $\mathbf{p}_k$ 是\emph{几何分量}（体素位置、三维点或网格面），$\boldsymbol\psi_k\in\mathbb{R}^{D}$ 是\emph{语义分量}——一个特征向量。把这一式与 \cref{ch:semantic} 的闭集元素 $m_k=(\mathbf{p}_k,\sigma_k)$ 并排看，差别一目了然：\textbf{只是把离散标签 $\sigma_k$ 换成了连续特征 $\boldsymbol\psi_k$}。其余一切——几何 $\mathbf{p}_k$ 由上游 SLAM（\cref{ch:vo,ch:lidar_slam,ch:dense}）给出、跨帧如何融合、如何查询——都顺着这一替换自然展开。本章余下各节，正是把这条主线在不同几何粒度（稠密 / 物体 / 隐式）与不同目标（最大表达力 / 最省内存 / 任务最优）上铺开。

\begin{insight}[几何与语义解耦：开放词汇建图为何能站在经典 SLAM 肩上]
\cref{eq:ow-mapelem} 藏着本章与全书衔接的关键：$\mathbf{p}_k$ 与 $\boldsymbol\psi_k$ \emph{解耦}。几何 $\mathbf{p}_k$ 仍由本书前面各章的 SLAM 后端解出（这也是为什么本章大多数方法严格说是\emph{纯建图}——它们假设已有上游 SLAM 提供带位姿的图像帧）；本章只负责往这副几何骨架上\emph{嵌一层特征}。正因如此，开放世界建图不必另起炉灶重做几何，而是\emph{接续}经典 SLAM 的输出——这与 \cref{ch:learning} 立的“学习只换母方程的零件、不推翻母方程”一脉相承：这里被“学习”接管的，是地图元素的\emph{语义分量}，而非几何与位姿。
\end{insight}

% ════════════════════════════════════════════════════════════════
\section{读懂开放词汇所需的基础模型（最小必要背景）}
\label{sec:openworld-fm}

\paragraph{动机：只取支点，不写综述。}
开放词汇建图的全部魔力，来自一类被称为\textbf{基础模型}（foundation model）的机器学习模型——在\emph{足够庞大}的数据集上训练、因而能合理期望它泛化到广泛工况的模型\cite{paull2026openworld}。完整的深度学习发展史有无数专著可循；本节先用 \cref{sec:openworld-fmhistory} 一段\emph{简史}把这条主线（AlexNet $\to$ ResNet $\to$ BERT/Transformer $\to$ 对比学习 $\to$ ViT $\to$ DINO $\to$ CLIP）勾出来——它回答“CLIP 这种图文对齐模型是怎么被一步步逼出来的”，是读懂后文的来龙去脉；随后只取理解后文所必需的\emph{三块支点}：(1) 什么是特征、为什么余弦相似度能度量“像不像”；(2) CLIP 如何把\emph{图像与文本}对齐到同一空间（后文一切开放词汇查询的支点）；(3) 类无关分割 SAM（后文从图像抠出物体的公共底座）。其余概念（LLM 的训练细节等）点到为止，需要深究者请查阅深度学习专著\cite{goodfellow2016deeplearning}。

\subsection{基础模型简史：从有监督单模态到自监督跨模态对齐}
\label{sec:openworld-fmhistory}

这段简史只为回答一个问题：\emph{CLIP 这种“图文同空间”的模型是怎么一步步被逼出来的}？读懂这条线，后文“把特征嵌进地图”就不再像凭空降下的黑魔法。全程参照本章蓝本的梳理\cite{paull2026openworld}。

\paragraph{第一步：有监督单模态（AlexNet $\to$ ResNet）。}
训练一个“基础模型”需要两样东西：能高效更新海量参数的\emph{深度网络架构}，与一个\emph{足够大}的数据集。计算机视觉里第一个勉强算得上的（虽通常不称作基础模型），是在 ImageNet（$320$ 万张\emph{人工标注}图）上有监督训练的卷积网络 AlexNet——它要靠人逐张贴标签。可纯卷积网络一深就犯\emph{梯度消失}：梯度穿过许多层后变得极小，靠它更新参数效率极低。\textbf{残差网络}（ResNet）\cite{he2016resnet}用一条\emph{跨层直连}缓解了这一点，让深网训练稳定高效；在 ImageNet 上预训练的 ResNet 视觉编码器，撑起了 YOLO、R-CNN 等一代\emph{闭集}视觉检测器（它们只认训练词典里的类别——正是本章要突破的对象）。

\paragraph{第二步：自监督让标注自动化（BERT $+$ Transformer）。}
人工标注 ImageNet 是苦役。自然语言处理（NLP）领域率先把它\emph{自动化}：\textbf{BERT}（Bidirectional Encoder Representations from Transformers）从语料里\emph{自造}带标签的任务——把句中一个词抠掉、让模型预测它（\emph{掩码语言建模}），于是英文维基（约 $25$ 亿词）与 BookCorpus（约 $8$ 亿词）无需任何人工标签就能拿来训练。这类\emph{无须人工标注}的训练，本书统称\emph{无监督/自监督}。BERT 还带火了一种新架构——\textbf{Transformer}，其核心是\emph{注意力}（attention）机制，能捕捉输入里\emph{长程}的依赖关系；与 BERT 几乎同时，第一代 \textbf{GPT}（Generative Pre-trained Transformer）也基于 Transformer、在 BookCorpus 上训练，显出相似能力。

\paragraph{第三步：把自监督搬上图像（对比学习 $+$ ViT $\to$ DINO）。}
语言的自监督直观（词有明确语义），但图像不然——\emph{一张图里随机抠出的小块未必有确定语义}。突破口是\emph{对比学习}：以 \textbf{SimCLR} 为例，把同一张图做两种\emph{数据增广}（抠块、翻转、变色\dots）得到两个视图，损失函数\emph{拉近}这两个增广视图（正样本）、\emph{推远}与其他图的视图（负样本）；如此学到的表示对增广\emph{不变}，因而利于下游视觉任务。与此并行，研究者把 Transformer 搬进视觉——\textbf{视觉 Transformer}（ViT，本书 \cref{sec:learning-depth} 已作单图深度网络主干用到它）把图像\emph{切成固定大小的图块}（tokenize）、配位置编码送进 Transformer，再用多头自注意力做分类，性能即可与最强的有监督 ResNet 比肩。两条线一合流，便有了 \textbf{DINO}\cite{caron2021dino}：用\emph{师生自蒸馏}（self-distillation，而非对比损失）无监督地训练一个 ViT 编码器，学出的视觉表示出奇地强。

\paragraph{第四步：跨模态对齐（CLIP）。}
对比学习还能跨\emph{模态}对齐表示空间。\textbf{CLIP}\cite{radford2021clip} 用从互联网收集的\emph{图文对}做对比训练，把一张图与它的文字说明在特征空间里\emph{拉近}、把与其他样本\emph{推远}——于是图像与文本被放进\emph{同一个}对齐的特征空间。这一步正是后文一切“以文查图”的开放词汇查询的支点（\cref{sec:openworld-clip} 详述）。至此，构建\emph{特征型}与\emph{生成式}基础模型所需的理论件已基本备齐；此后的进步主要来自模型、数据、算力的持续放大，是否能维持同样的进步速率尚不明朗。

\begin{note}[一条要记住的暗线：互联网级数据是命根]
这段简史里反复出现的关键，不是某个巧妙的网络，而是\emph{互联网级数据}——正是它喂出了这些模型惊人的表征力。这对 SLAM 是个\emph{警示}：本书前面各章倚重的若干模态（激光、雷达等）\emph{并没有}如此唾手可得的海量数据\cite{paull2026openworld}。这也部分解释了为何开放词汇建图\emph{先}在“图像 $+$ 文本”这一数据最富集的模态上开花。本章用到的 ViT、RAFT 等学习件与本部 \cref{ch:learning} 同源（\cref{sec:learning-flow} 的 RAFT、\cref{sec:learning-depth} 的 ViT），读者可对照着看“学习如何接进几何”。
\end{note}

\subsection{特征、嵌入与余弦相似度}
\label{sec:openworld-feature}

所谓\textbf{特征}（feature，亦称嵌入 embedding、表示 representation、编码 code），是一个深度神经网络对输入数据点编码后得到的\emph{实值向量}\cite{paull2026openworld}。形式化地：对模态 $X$ 的输入 $x$，编码器 $\mathcal{E}_X$（通常是一个深度网络）给出 $D$ 维特征
\begin{equation}
  \boldsymbol\psi_x=\mathcal{E}_X(x)\in\mathbb{R}^{D}.
  \label{eq:ow-encoder}
\end{equation}
对本章最要紧的，是这类特征的一条性质：\emph{特征空间里的距离能度量数据点之间有意义的相似度}。这让我们能回答“这两段文字有多像？”“这句话和这张图有多像？”之类的问题。最常用的度量是\textbf{余弦相似度}——把两个特征看成从原点出发的向量、取其夹角的余弦：若 $\boldsymbol\psi_a,\boldsymbol\psi_b$ 已单位化（$\|\boldsymbol\psi\|=1$，下同），则余弦相似度就是内积 $\langle\boldsymbol\psi_a,\boldsymbol\psi_b\rangle$。本章一律假设特征已单位化，故以内积表余弦。

\begin{note}[为什么是余弦而非欧氏距离]
开放词汇特征里，“概念是否相同”更多体现在\emph{方向}而非\emph{长度}上：同一概念的不同实例（远近、明暗不同的同一把椅子）其特征向量长度可能不同，但方向相近。余弦相似度只看方向、对长度不敏感，故成为特征比对的默认选择。这也是为什么各模型常把输出特征\emph{单位化}——把比较统一到单位球面上。
\end{note}

\subsection{多模态对齐：CLIP 的共享嵌入空间}
\label{sec:openworld-clip}

单一模态内的特征已能比相似度，但开放词汇查询要的是\emph{跨模态}——用一句\emph{文本}去查一张\emph{图}（或一片三维点云）。这要求图像与文本被映到\emph{同一个}对齐的特征空间。\textbf{CLIP}（Contrastive Language--Image Pre-training）是这一思路在视觉–语言上的标杆\cite{radford2021clip}。它训练两个编码器——图像编码器 $\mathcal{E}_{\mathrm{img}}$ 与文本编码器 $\mathcal{E}_{\mathrm{txt}}$——并用一个\emph{对比损失}，把\emph{来自同一图文对}的图像特征与文本特征在空间里\emph{拉近}（正样本），把图像特征与\emph{其他随机文本}的特征\emph{推远}（负样本）。在从互联网收集的四亿图文对上这样训练后，得到的特征能开箱表征极广的视觉与文本概念\cite{radford2021clip}。

关键在于这套训练\emph{不依赖任何预定义的类别集}。于是给定一张图 $I$ 与一段文本 $T$，CLIP 抽出单位化特征 $\boldsymbol\psi_I=\mathcal{E}_{\mathrm{img}}(I)$、$\boldsymbol\psi_T=\mathcal{E}_{\mathrm{txt}}(T)$，二者的余弦相似度 $\langle\boldsymbol\psi_I,\boldsymbol\psi_T\rangle$ 就度量了“这张图和这句话有多匹配”。这立刻给出一个\emph{零样本分类器}：把一组类别标签写成文本形式 $\mathsf{C}$（可套模板，如“一张{狗}的照片”），则对图像 $I$ 的预测类别为
\begin{equation}
  \hat{c}=\operatorname*{arg\,max}_{c\in\mathsf{C}}\ \langle\boldsymbol\psi_I,\ \boldsymbol\psi_c\rangle,\qquad \boldsymbol\psi_c=\mathcal{E}_{\mathrm{txt}}(c).
  \label{eq:ow-clip-cls}
\end{equation}
要把分类器改用到新问题，只需\emph{重新定义文本提示集} $\mathsf{C}$，无需任何重训。\cref{eq:ow-clip-cls} 是本章后续所有“零样本查询”的母型——把图 $I$ 换成地图元素的特征 $\boldsymbol\psi_k$，就得到 \cref{sec:openworld-mapping} 的地图查询。

\begin{insight}[CLIP 共享空间：本章一切查询的支点]
请把 \cref{eq:ow-clip-cls} 牢牢记住：它说\textbf{图像（或任何有对齐编码器的模态）与文本可以被放到同一把尺上比相似度}。这就是开放词汇建图的全部支点——把地图元素的特征 $\boldsymbol\psi_k$ 嵌成 CLIP 空间里的向量后，任何一句话也能编码成同一空间的向量，于是“地图里哪个元素最匹配这句话”就化成一次余弦近邻搜索（\cref{eq:ow-query}）。$\mathsf{C}$ 不必在建图时已知、可随时重定义——这正是\emph{开集}相对\emph{闭集}的根本自由。视觉-only 的特征（如 DINO，自蒸馏训练的视觉编码器）也很强，但因\emph{不与语言对齐}，只支持“以图查图 / 点击查询”，不支持“以文查图”——这正凸显了 CLIP 的图文对齐为何不可或缺。
\end{insight}

\paragraph{把 CLIP 放回特征模型家族：文本侧。}
CLIP 之所以特殊，在于它\emph{跨模态对齐}；但它两侧的“零件”——纯文本与纯视觉的特征模型——本身也是后文常用的工具，且恰好衬出 CLIP 对齐的不可替代。先看\emph{文本侧}：\cref{sec:openworld-fmhistory} 提到的 BERT 给出的是\emph{词元级}（token，即输入的子串）特征，要比较\emph{整句}的相似度并不直接。\textbf{SentenceBERT} 在 BERT 上加一个\emph{池化}操作并微调，使其直接输出\emph{句级}特征——两句话的相似度于是可由它们句向量的\emph{余弦相似度}（\cref{sec:openworld-feature}）高效算出\cite{paull2026openworld}。SentenceBERT 有多语种、多任务的现成权重。这正是 \cref{sec:openworld-objgraph} 用“物体的自然语言描述”做检索时，背后“句子比句子”的度量来源。

\paragraph{视觉侧：DINO 的图像级与 patch 级特征。}
再看\emph{视觉侧}：\cref{sec:openworld-fmhistory} 的 DINO\cite{caron2021dino} 用师生自蒸馏训练 ViT，其训练的一味关键作料是\emph{多裁剪增广}（multi-crop）——教师只看\emph{全局视图}（$>50\%$ 的图），学生既要处理\emph{局部视图}（$<50\%$）也要处理全局视图，逼出“\emph{局部到全局}”的一致表示。DINO 默认输出\emph{图像级}视觉特征，零样本就能支撑图像分类、图像检索（特征空间里的距离可直接拿来做 $k$ 近邻分类器，甚至逻辑回归都表现不俗）。\textbf{更要紧的是}：撑起 DINO 的 Transformer 架构还能输出 \emph{patch 级}特征（如对每个 $8\times8$ 图块给一个向量）；在这些\emph{稠密}特征上再训小模型，可在\emph{像素级}下游任务——\emph{语义分割}、\emph{深度估计}、\emph{表面法向估计}——上取得强性能\cite{oquab2024dinov2}。\textbf{DINOv2}\cite{oquab2024dinov2} 进一步改良训练、并整理出一个 $1.42$ 亿张图的精选数据集，得到一个对图像有强\emph{整体理解}的视觉特征提取器。

\begin{insight}[为何 DINO 强却仍离不开 CLIP：对齐 vs 不对齐]
DINO 的视觉特征非常强，但它\emph{只在视觉模态内}训练、\textbf{不与语言对齐}——故只支持“以图查图 / 点击查询”，\emph{不}支持“以文查图”。这里还有一个值得记住的反差：从 CLIP 的视觉编码器\emph{也}能抽\emph{稠密 patch 级}特征，但\emph{不应指望它们语言对齐}，且有研究发现它们在视觉与三维相关任务上不如 DINO\cite{oquab2024dinov2}。一句话：\emph{要“以文查图”，CLIP 的图文对齐不可替代；要纯视觉的稠密 patch 语义，DINO 往往更好}。这解释了为何后文（如 \cref{sec:openworld-implicit} 的 LERF 用 DINO 监督使物体边界更清晰）常\emph{两者并用}——CLIP 管语言对齐、DINO 管边界与几何。
\end{insight}

\paragraph{生成式基础模型：一句话带过。}
与上述\emph{特征型}模型并列的，是\emph{生成式}基础模型：大语言模型（LLM）以“预测下一个词”为目标在海量文本上训练，能续写、对话、生成代码、做符号规划；视觉-语言模型（VLM，如把 CLIP 视觉编码器接到 LLM 上的 LLaVA）则额外接受图像输入、能描述图像、做视觉问答\cite{paull2026openworld}。本章在两处会用到它们：一是 \cref{sec:openworld-objgraph} 用 VLM 为物体生成\emph{自然语言描述}、用 LLM 做\emph{基于描述的检索}（以处理含否定、可供性的复杂查询）；二是 \cref{sec:openworld-future} 讨论“把地图喂给 LLM”。它们的内部机理（注意力、自回归生成）非本章重点，按下不表。

\subsection{类无关分割：SAM}
\label{sec:openworld-sam}

后文几乎所有方法都需要先把图像\emph{抠成若干区域}，再对每个区域抽特征——因为基础模型直接吃整张图，给出的是“整幅画里所有物体的混合”特征，而我们要的是\emph{逐物体、乃至逐像素}的细粒度语义。完成这一步的工具是\textbf{类无关分割}（class-agnostic segmentation）：检出一组掩膜 $\mathrm{Seg}(I)=\{Z_i\}_{i=1}^{R}$、$Z_i\in\{0,1\}^{H\times W}$，每个掩膜圈出图中一个可能的物体。关键是\emph{类无关}——不预设物体属于哪一类，才能不漏掉机器人环境里千奇百怪的物体、也才配得上前面基础模型的开放词汇本性\cite{paull2026openworld}。

实现这一函数的主力是 \textbf{SAM}（Segment Anything Model）\cite{kirillov2023sam}：在大规模监督数据上训练的\emph{可提示}分割模型，能接受像素坐标、边界框乃至另一个掩膜作为提示、输出对应掩膜，也能自动处理整幅图、输出描述各实例的一组掩膜。SAM 既可独立使用，也可串在一个开放词汇检测器之后（用检测框作 SAM 的提示），以获得更好的物体边界并提速。本章一律把这一步记作 $\mathrm{Seg}(\cdot)$，视其为给定的黑盒。

\begin{insight}[三块支点，一条数据流]
本节三块连起来就是后文每个系统的\emph{共同前半段}：$\mathrm{Seg}(\cdot)$（SAM）把图抠成区域 $\to$ $\mathcal{E}_{\mathrm{img}}$（CLIP）对每个区域抽\emph{开放词汇特征} $\to$ 用针孔反投影（\cref{ch:camera}）把像素特征抬进三维。后文稠密、物体、隐式三类方法，分歧只在“特征存到哪一级几何上、跨帧怎么融、怎么查”——前半段几乎一致。记住这条数据流，下一节就不会迷失在系统名里。
\end{insight}

% ════════════════════════════════════════════════════════════════
\section{开放词汇建图}
\label{sec:openworld-mapping}

\paragraph{动机：把特征织进地图。}
本节是全章技术主体。设定与 \cref{ch:semantic} 平行而更进一步：主传感器为 RGB-D 相机（既能由图像抽特征、又能逐像素抬进三维），假设有一串\emph{带位姿}的 RGB-D 帧（位姿来自上游 SLAM）、场景\emph{静态}，逐帧增量建图\cite{paull2026openworld}。我们要回答三件事，对应一条流水线的三段：(i) 用什么三维表示存特征；(ii) 怎么从一帧抽出开放词汇特征；(iii) 跨帧怎么把对应同一三维位置的特征\emph{融}成一个、以及怎么\emph{查}。先在稠密表示上把这三段讲透（\cref{sec:openworld-dense}），再看物体级 / 场景图（\cref{sec:openworld-objgraph}）与隐式（\cref{sec:openworld-implicit}）如何在同一内核上变体。

\subsection{稠密开放词汇地图}
\label{sec:openworld-dense}

\paragraph{(i) 表示：往稠密几何里塞特征。}
稠密开放词汇地图，就是把基础模型的语义特征嵌进 \cref{ch:dense} 介绍过的某种稠密表示——无序三维点集、TSDF 体素网格、或投影到二维栅格。选哪种由下游任务决定：导航偏好二维栅格的逐格聚合，移动操作偏好三维体素\cite{paull2026openworld}。无论哪种，地图元素都形如 \cref{eq:ow-mapelem} 的 $m_k=(\mathbf{p}_k,\boldsymbol\psi_k)$，且\emph{几何分量 $\mathbf{p}_k$ 由上游 SLAM 提供}（\cref{sec:openworld-closed-open} 已强调的解耦）。

\paragraph{(ii) 特征抽取：全局--局部融合。}
一个直接做法是用\emph{直接输出逐像素特征}的模型（在分割数据上微调过的 CLIP 变体）。但近来的证据表明，这类经微调的稠密特征对\emph{长尾}概念的覆盖不如原始未改动的 CLIP 特征\cite{paull2026openworld}。于是另一条路（以 ConceptFusion 为代表\cite{jatavallabhula2023conceptfusion}）选择\emph{保留}原始 CLIP 特征、用\emph{特征算术}拼出逐像素特征，以平衡“全局语境”与“局部细节”两种需求：

\begin{itemize}
  \item \textbf{全局特征} $\boldsymbol\psi^{G}=\mathcal{E}_{\mathrm{img}}(I)$：把整幅图编码，给出全图语境；
  \item \textbf{局部特征} $\boldsymbol\psi_i^{L}=\mathcal{E}_{\mathrm{img}}(I_i)$：对 $\mathrm{Seg}(I)$ 得到的每个掩膜，转成边界框、裁出紧致图块 $I_i$ 再编码，给出局部细节；
  \item \textbf{逐掩膜混合} $\boldsymbol\psi_i^{P}=w_i\,\boldsymbol\psi^{G}+(1-w_i)\,\boldsymbol\psi_i^{L}$：权重 $w_i\in[0,1]$ 这样选——\emph{凸显那些既不同于全局、又区别于其他局部的掩膜}（即该局部越“独特”，越偏重它自己的局部特征）；
  \item \textbf{逐像素特征} $\boldsymbol\psi_{u,v}^{P}$：把覆盖像素 $(u,v)$ 的所有掩膜的混合特征相加（即所有满足 $Z_i[u,v]=1$ 的 $\boldsymbol\psi_i^{P}$），再单位化。
\end{itemize}

\noindent 配上针孔模型与逐像素的特征-坐标对应，就按 \cref{sec:dense-rgbd-pc} 的流程生成一帧\emph{带开放词汇特征的有序点云}：每个三维点配一个语义特征。仿照 \cref{ch:semantic} 的 \emph{range-category 观测}，我们把这种“每个测距点附带一个特征向量”的观测称作 \textbf{range-feature 观测}\cite{paull2026openworld}。这一命名不是修辞——它点明了本章与上一章的\emph{平行结构}：上一章测距点附带的是\emph{类别} $\sigma$，本章附带的是\emph{特征} $\boldsymbol\psi$。

\begin{insight}[range-category $\to$ range-feature：同一脉络的连续化]
\cref{ch:semantic} 的多类贝叶斯融合，把“每个测距点带一个类别”的 range-category 观测，递归融成逐元素的\emph{类别分布}（多类 log-odds，\cref{eq:dense-octo-logodds} 的多类推广）。本章把类别换成连续特征，于是 range-category 变 range-feature、类别分布变\emph{特征向量}。这条“逐元素证据累积”主线——从二值占据（\cref{eq:dense-octo-logodds}）到多类分布（\cref{ch:semantic}）到连续特征（本节）——是同一种智慧在越来越富表达的语义上的三次落地。读者若把它看成三套互不相干的方法，便错失了贯穿三章的统一性。
\end{insight}

\paragraph{(iii) 融合与查询。}
跨帧融合要解决：不同帧里落到\emph{同一}三维位置的多个特征，怎么合成一个标注该地图元素的特征。最朴素的做法（如 VLMaps）是把对应同一二维栅格的特征\emph{直接平均}\cite{paull2026openworld}。ConceptFusion 用一个\emph{置信加权}的递归平均：把 range-feature 观测转到全局系后，对每个地图元素 $m_k$，用置信度 $c_k$ 把历史特征与新观测 $\boldsymbol\psi_{u,v}^{P}$ 加权融合\cite{jatavallabhula2023conceptfusion}：
\begin{equation}
  \boldsymbol\psi_k\ \leftarrow\ \frac{c_k\,\boldsymbol\psi_k+\alpha\,\boldsymbol\psi_{u,v}^{P}}{c_k+\alpha},\qquad
  c_k\ \leftarrow\ c_k+\alpha,\qquad
  \alpha=\exp\!\Big(\!-\frac{\gamma^2}{2\sigma^2}\Big),
  \label{eq:ow-fusion}
\end{equation}
其中 $\alpha$ 是每个像素-特征的置信度，$\gamma$ 是到相机中心的径向距离、$\sigma$ 是尺度参数。直觉：\emph{越靠近相机光轴中心的特征越可信}（径向越远、畸变与采样越差），与三维重建里的经验一致。

\begin{insight}[\cref{eq:ow-fusion} 就是“信息加权平均”，与本书融合主线同根]
\cref{eq:ow-fusion} 不必当成新公式。把 $c_k$ 读作“已累积的置信权重”、$\alpha$ 读作“本次观测的权重”，它就是一个\emph{权重递增的加权平均}——与 \cref{ch:dense} 的 TSDF 加权融合（按权重累加、$W_\eta$ 封顶）、乃至卡尔曼更新“按信息量加权融合先验与观测”是同一种智慧：\textbf{可信的观测多说话，不可信的少说话}。差别只在融合的对象从“距离/占据”换成了“高维特征向量”。这也解释了为什么有方法（见下）宁可不平均高维向量——平均会把语义“抹糊”。
\end{insight}

平均高维特征有两个隐患：一是把不同概念的特征\emph{平均}会得到一个语义模糊的向量，二是对离群特征不鲁棒。于是有方法改用\emph{非平均}的聚合：为每个地图元素维护一个特征\emph{集合} $\{\boldsymbol\psi_{k,1},\dots,\boldsymbol\psi_{k,t}\}$，取\emph{中位元}（medoid，到集合内其余特征欧氏距离之和最小者）作代表；或用聚类取若干代表中心、或选最接近多数簇质心者以增鲁棒\cite{paull2026openworld}。这些都在“逐元素聚合证据”的同一框架内，只是把“加权平均”换成了对离群更稳的统计量。

\textbf{查询}是开放词汇地图的主用例，也是 \cref{eq:ow-clip-cls} 在地图上的直接落地。两类查询：

\emph{零样本分割}：把候选类别写成文本提示集 $\mathsf{C}$、编码成 $\boldsymbol\psi_C=\mathcal{E}_{\mathrm{txt}}(C)$，则地图元素 $m_k$ 的预测标签
\begin{equation}
  \hat{y}_k=\operatorname*{arg\,max}_{C\in\mathsf{C}}\ \langle\boldsymbol\psi_k,\ \boldsymbol\psi_C\rangle.
  \label{eq:ow-query}
\end{equation}
这与 \cref{eq:ow-clip-cls} 的图像分类器\emph{完全同形}——只是把整图特征 $\boldsymbol\psi_I$ 换成了地图元素特征 $\boldsymbol\psi_k$。\textbf{关键在于 $\mathsf{C}$ 不必在建图时已知、可随时重定义}以预测不同标签——这正是“开放世界”的全部自由。

\emph{开放式查询}：把一段\emph{自由文本}查询 $T$ 编码成 $\boldsymbol\psi_T=\mathcal{E}_{\mathrm{txt}}(T)$，用相似度 $\langle\boldsymbol\psi_k,\boldsymbol\psi_T\rangle$ 给所有地图元素打分、取阈值返回相关元素（其空间坐标还可聚类以分出不同实例）。查询不限于文本：任何有对齐编码器的模态都行——\emph{点击查询}（取地图中被点元素的特征）、\emph{以图查询}（$\boldsymbol\psi_{I}=\mathcal{E}_{\mathrm{img}}(I_{\mathrm{query}})$）、乃至\emph{以声查询}（用音频编码器）。这正是本章开篇图景“用一句‘适合喝咖啡看书的地方’查到所有沙发和椅子”背后的机制。

\subsection{物体级地图与场景图}
\label{sec:openworld-objgraph}

\paragraph{动机：稠密太贵，按物体存。}
稠密开放词汇地图能在极细粒度上建模语义，但\emph{随环境增大而内存爆炸}——高维特征逐点/逐体素地存，代价高昂\cite{paull2026openworld}。然而注意一个事实：图像里\emph{相邻像素往往属于同一物体、共享语义}（这正是前面用掩膜抽局部特征的依据）；推广到地图——\emph{相邻地图元素往往语义相近}。于是可以把特征存到\emph{物体}这一级，\emph{一物一向量}，内存骤降，且天然契合具身 AI 的\emph{物体检索}与\emph{以物体为中心的任务规划}。

\paragraph{表示与流水线。}
把地图表示成 $J$ 个物体的集合 $\mathcal{V}_O=\{o_j\}_{j=1}^{J}$，每个物体 $o_j$ 配一个几何表示（物体中心、三维框、或稠密点云 $\mathsf{P}_{o_j}$）、一个特征 $\boldsymbol\psi_{o_j}$、一个检测计数 $n_{o_j}$\cite{paull2026openworld}。增量流水线（以 ConceptGraphs\cite{gu2024conceptgraphs}、OVIR-3D 为例）每来一帧：

\begin{enumerate}
  \item \textbf{抽掩膜特征}：$\mathrm{Seg}(I)$ 得掩膜，转框、裁块、$\mathcal{E}_{\mathrm{img}}$ 抽掩膜特征 $\boldsymbol\psi_i^{M}$（与 ConceptFusion 抽局部特征同法）；用深度、位姿、内参把每个掩膜的像素反投影到三维，得掩膜点云 $\mathsf{P}_i^{M}$（可用 DBSCAN 去噪）。于是得到\emph{掩膜-特征观测} $\{(\mathsf{P}_i^{M},\boldsymbol\psi_i^{M})\}$。
  \item \textbf{物体关联}：衡量新掩膜与已有物体的相似度。\emph{几何相似度} $\phi_{\mathrm{geo}}(\mathsf{P}_i^{M},\mathsf{P}_{o_j})\in[0,1]$ 取“$\mathsf{P}_i^{M}$ 中在阈值 $\delta_{\mathrm{nn}}$ 内于 $\mathsf{P}_{o_j}$ 有最近邻的点所占比例”；\emph{语义相似度}由特征相似度的函数给出，例如 $\phi_{\mathrm{sim}}(\boldsymbol\psi_i,\boldsymbol\psi_{o_j})=\tfrac12(\langle\boldsymbol\psi_i,\boldsymbol\psi_{o_j}\rangle+1)\in[0,1]$（把 $[-1,1]$ 的余弦映到 $[0,1]$）。两者合成总分或分别取阈，判定掩膜-物体匹配；无匹配则用该掩膜初始化\emph{新物体}。
  \item \textbf{物体融合}：若第 $i$ 个掩膜匹配上 $o_j$，则按计数加权融合特征、并入点云、计数加一：
  \begin{equation}
    \boldsymbol\psi_{o_j}\leftarrow\frac{n_{o_j}\,\boldsymbol\psi_{o_j}+\boldsymbol\psi_i^{M}}{n_{o_j}+1},\qquad
    \mathsf{P}_{o_j}\leftarrow\mathsf{P}_{o_j}\cup\mathsf{P}_i^{M},\qquad
    n_{o_j}\leftarrow n_{o_j}+1,
    \label{eq:ow-objfuse}
  \end{equation}
  点云随后降采样去冗余。这与稠密融合 \cref{eq:ow-fusion} 同构，只是权重换成检测计数。
  \item \textbf{后处理}：每隔若干帧在 $\mathcal{V}_O$ 内部跑关联-融合以并掉冗余物体；删掉检测计数 $n_{o_j}$ 过低的物体（多为噪声）。
\end{enumerate}

\noindent 查询机制与稠密情形几乎一致：零样本分割时三维点继承所属物体的标签 $\hat{y}_k=\arg\max_{C}\langle\boldsymbol\psi_{o_j},\boldsymbol\psi_C\rangle$；\emph{物体检索}则更有用——给定文本查询 $T$，取最相关物体 $o^\star=\arg\max_{o}\langle\boldsymbol\psi_o,\boldsymbol\psi_T\rangle$，其点云可直接交给运动规划去导航、检视、操作。“找点能喝的”就会与一罐汽水的特征强相关。

\paragraph{特征之外：用语言描述物体。}
语言对齐的物体特征，在面对\emph{短的描述性}查询时最灵；但对含\emph{否定}或\emph{可供性}的复杂查询常失灵——“一罐汽水\emph{之外}能喝的东西”很可能仍与汽水特征强相关，因为特征不“懂”否定\cite{paull2026openworld}。一条出路是\emph{直接用自由语言描述物体}：借 VLM 之力，跟踪每个物体贡献点最多的若干图块、连同提示“描述图中央的物体”喂给 VLM，逐视图描述后再用 LLM 汇总成每物体一句描述。物体描述远比闭集标签富表达；查询时让 LLM 在物体描述列表里“检索”最相关者，能处理特征式检索搞不定的复杂查询（代价是 VLM/LLM 推理比特征比对慢得多）。有方法\emph{同时}用特征与语言描述，各取所长。

\paragraph{物体场景图与分层场景图。}
更进一步，把地图表示成图 $\mathcal{G}=(\mathcal{V}_O,\mathcal{E}_O)$，用边集 $\mathcal{E}_O$ 显式建模物体间\emph{关系}，即三维\emph{场景图}（与 \cref{ch:semantic} 引入的场景图同一记号、同一思想）。开放词汇下，边的描述也可开集化：ConceptGraphs 用物体描述与位置、让 LLM 推断“$a$ 在 $b$ 上”“$b$ 在 $a$ 里”之类空间关系，并可推广到“背包可能放在衣柜里”这类语言模型能解释的关系\cite{gu2024conceptgraphs}；边特征也可由“两物体共现图像的编码”给出，或用图神经网络学习关系。

正交于物体间关系的，是把场景\emph{分层划分}（沿用 \cref{ch:semantic} 的分层图定义）。以 HOV-SG\cite{werby2024hovsg} 为代表，把室内环境分成\emph{楼层--区域--物体--可导航自由空间}多层，节点与边按层组织：
\begin{equation}
  \mathcal{V}=\mathcal{V}_E\cup\mathcal{V}_F\cup\mathcal{V}_R\cup\mathcal{V}_O,\qquad
  \mathcal{E}=\mathcal{E}_{EF}\cup\mathcal{E}_{FR}\cup\mathcal{E}_{RO},
  \label{eq:ow-hsg}
\end{equation}
其中 $\mathcal{V}_E$ 是根节点，$\mathcal{V}_F,\mathcal{V}_R,\mathcal{V}_O$ 分别是楼层、区域、物体节点，边按“根--楼层--区域--物体”逐层相连。分层带来更小的内存、更快的图搜索、以及\emph{逐层语言定位}：每个概念节点都配至少一个由 CLIP 给出的开放词汇特征（楼层套模板“第{X}层”，区域用若干代表视图的 CLIP 特征经 k-means 蒸馏成 $K$ 个代表、再对候选区域类别集 $\mathsf{C}_R$ 打分投票得区域类别，如 \cref{eq:ow-clip-cls} 之于区域）。于是面对“去一楼客厅里的那盆植物”这样的长查询，可由 LLM 先拆成“植物 / 客厅 / 一楼”几个概念，再\emph{自顶向下}（楼层 $\to$ 区域 $\to$ 物体）逐层用 CLIP 相似度匹配——把一个复杂空间查询分解为几次 \cref{eq:ow-query} 式的逐层检索。

\paragraph{旁支：把语义挂到视点上（Pose Maps）。}
本小节到此都把开放词汇特征挂在\emph{场景几何}上（稠密的逐点/逐体素、物体的一物一向量、场景图的分层节点）。还有一类相关方法走另一条更稀疏的路——\textbf{位姿 / 视点级地图}（Pose Maps）：干脆把开放词汇语义\emph{挂到相机位姿（视点）上}，每个关键帧 / 视点存一个（或一组）由 $\mathcal{E}_{\mathrm{img}}$ 抽出的特征，地图便是\emph{一串带语义的视点}\cite{paull2026openworld}。这种表示比物体级更\emph{稀疏}，且天然接住本书已有的位姿图机器：它可直接用于\emph{房间分割}（按视点特征聚类划区域）、\emph{回环检测}（\cref{ch:loop}：视点特征相近即疑似重访同处）与\emph{重定位}。查询时它不返回“哪个三维点 / 物体最匹配”，而返回\emph{最相关的视点}——这恰好可拿来\emph{指定导航目标}（“去最像‘厨房水槽’的那个视点”）。可见“特征嵌进地图”这条内核，连\emph{挂在哪一级几何}都能一路稀疏化：逐点 $\to$ 一物一向量 $\to$ 一帧一向量，仍是同一道“表达力换内存 / 可检索性”的权衡。

\begin{insight}[物体级与稠密：同一内核，不同粒度]
别被 ConceptFusion / ConceptGraphs / HOV-SG 等系统名晃花了眼——它们共享同一内核：$\mathrm{Seg}\to$ 抽特征 $\to$ 反投影 $\to$ 关联融合 $\to$ 余弦查询。分歧只在\emph{特征存到哪一级几何}：稠密（逐点/逐体素，最富表达、最贵）、物体（一物一向量，省内存、利检索）、分层场景图（物体再往上聚成区域/楼层，利大尺度语言导航）。\textbf{粒度自低向高，是用表达力换内存与可检索性。}这恰好引出下一节的根本问题——粒度到底该取多少，谁说了算？
\end{insight}

\subsection{隐式表示（选读）}
\label{sec:openworld-implicit}

\cref{ch:dense} 末尾介绍过用神经辐射场（NeRF）与三维高斯泼溅（3DGS）建可微渲染地图（其归属本书在结语 \cref{ch:epilogue} 指明）。这两类表示也被扩展以容纳开集语义：NeRF 系给隐式辐射场\emph{添一路语义特征输出}，3DGS 系给每个高斯\emph{挂一个特征}、再光栅化到像平面\cite{paull2026openworld}。由于给每个高斯存完整 CLIP 特征（如 ViT-L/14 的 $768$ 维）内存巨大，常见做法是\emph{学一个压缩的} CLIP 特征、或把高斯分组、每组共享一个特征向量。

以 LERF\cite{kerr2023lerf}（NeRF 系）与 LangSplat\cite{qin2024langsplat}（3DGS 系）为例。LERF 给 NeRF 加一路语义嵌入输出，输入为点 $\mathbf{x}$ 与尺度 $s$（以 $\mathbf{x}$ 为中心的立方体宽度），语义嵌入像 NeRF 的颜色一样沿光线\emph{体渲染}（去掉视角依赖以保证不同视角下输出一致）；训练时对每幅训练图在\emph{多个随机尺度}的图块上预计算 CLIP 特征构成“语义金字塔”作监督，并辅以 DINO 监督以使物体边界更清晰。查询时对渲染图中每个语义嵌入 $\boldsymbol\psi_{\mathrm{render}}$ 与文本嵌入 $\boldsymbol\psi_T$ 算一个\emph{相关性分数}
\begin{equation}
  \min_{i}\ \frac{\exp\!\big(\langle\boldsymbol\psi_{\mathrm{render}},\boldsymbol\psi_T\rangle\big)}
  {\exp\!\big(\langle\boldsymbol\psi_{\mathrm{render}},\boldsymbol\psi_{\mathrm{canon}}^{i}\rangle\big)+\exp\!\big(\langle\boldsymbol\psi_{\mathrm{render}},\boldsymbol\psi_T\rangle\big)},
  \label{eq:ow-lerf}
\end{equation}
其中 $\boldsymbol\psi_{\mathrm{canon}}^{i}$ 是一组\emph{规范短语}（如“object”“things”“stuff”“texture”）的特征，用以为相关性分数去噪——直觉是“查询词相对一组中性词有多突出”。LangSplat 则对 SAM 在\emph{三种粒度}下的分割块算 CLIP 特征（既抓局部又抓全局、且无需 LERF 那样另加 DINO 监督），并训一个\emph{场景专属}的自编码器把 CLIP 特征从 $512$ 维压到 $3$ 维存进高斯（直觉：单个场景的语义变化远小于 CLIP 训练语料，故低维足矣），查询时按三种粒度各算 \cref{eq:ow-lerf} 取分最高者。

\begin{insight}[隐式表示仍是同一内核的变体]
\cref{eq:ow-lerf} 的相关性分数本质仍是 \cref{eq:ow-query} 的余弦查询，只是加了一组规范短语作归一化、并在\emph{渲染}出的嵌入上做。隐式表示换掉的是\emph{几何载体}（从离散点/体素换成连续辐射场/高斯）与\emph{特征如何落到三维}（体渲染监督而非反投影累积），“嵌特征 $+$ 余弦查”这条内核没变。这与 \cref{ch:dense} 把 NeRF 体渲染与 OctoMap 对数几率认作“同根”（连续 vs 离散的占据渲染）是同一种看法：\emph{隐式只是把离散表示连续化}。
\end{insight}

% ════════════════════════════════════════════════════════════════
\section{任务驱动的表示：用信息瓶颈压缩地图}
\label{sec:openworld-task}

\paragraph{动机：场景到底该切成几个物体？}
开集建图比闭集能表达远更丰富的概念，却带来一个闭集时代不存在的新难题：\textbf{物体的粒度是病态的}\cite{paull2026openworld}。闭集时，物体由词典里的标签定义（“工具”“杯子”），粒度随词典定死。开集时，“正确的物体粒度”却没有先验答案。一架钢琴、一间空屋：要\emph{搬钢琴}，整架琴是\emph{一个}物体的粒度刚好；要\emph{弹钢琴}，场景须建成\emph{九十多个}物体——每个琴键一个；要\emph{调音}，又得是\emph{上百个}物体——每根弦、每个弦轴。同理，一片森林该建成一块地貌、还是一根根枝、一片片叶、一根根干？\textbf{在指定“表示要支持什么任务”之前，这个问题无解。}多数实践者都同意：正确的场景表示\emph{取决于智能体要完成的任务}——但在视觉-语言基础模型出现前，一直不清楚如何把“任务知识”切实纳入建图。

\begin{insight}[这正是“鲁棒感知时代”的“任务驱动感知”]
请回看 \cref{tab:intro-eras}：开篇引 Cadena 等划出的“鲁棒感知时代”四诉求，最后一条就是\emph{任务驱动的感知}（task-driven perception）——感知不该一味追求“重建一切”，而该\emph{按任务需要}决定建什么、建多细。本节正是这一诉求在开放世界地图上的\emph{数学落地}：把“场景该切多细”写成一个由任务驱动的\emph{优化问题}。本书走到这里，开篇那张“时代表”上最抽象的一格，终于有了可计算的形式。
\end{insight}

\paragraph{用信息瓶颈把粒度形式化。}
前面物体级地图（如 \cref{sec:openworld-objgraph} 的 ConceptGraphs）按语义/几何相似度聚类来定物体——这对某些场景够用，但不能\emph{一般地}保证粒度正确。Clio\cite{maggio2024clio} 把\textbf{任务驱动开集建图}正式化为：给定一组自然语言任务（如“清理背包”“取调料包”），构造一张物体粒度\emph{恰好支持这些任务}的地图。其关键洞见是——这个问题可以用经典的\textbf{信息瓶颈}（Information Bottleneck，IB）理论\cite{tishby1999ib}精确表述。

设有一组细粒度\emph{基元}观测 $\mathcal{X}$（极端情形下可以是所有视觉观测的全部像素，实践中是过分割的类无关掩膜对应的三维基元），目标是形成一个\emph{压缩}的场景表示 $\tilde{\mathcal{X}}$（即任务相关的物体），它\emph{既尽量压缩、又尽量保留完成任务 $\mathcal{Y}$ 所需的信息}。把“基元 $x$ 被聚到物体 $\tilde{x}$”的软分配写成概率 $p(\tilde{x}\mid x)$，则任务驱动聚类问题写成信息瓶颈：
\begin{equation}
  \min_{p(\tilde{x}\mid x)}\ \ I(\mathcal{X};\tilde{\mathcal{X}})\ -\ \beta\,I(\tilde{\mathcal{X}};\mathcal{Y}),
  \label{eq:ow-ib}
\end{equation}
其中 $I(\cdot;\cdot)$ 是\emph{互信息}，$\beta>0$ 调节“压缩”与“保任务信息”之间的权衡\cite{maggio2024clio}。

\begin{insight}[\cref{eq:ow-ib} 接住的，正是本书的互信息]
这条式子并非空降的新理论。本书早在状态估计章就用本书记号定义了互信息 $I(\mathbf{x},\mathbf{y})$（\cref{eq:mutual-info}），并用它刻画\emph{主动感知}——“一次观测能消去多少不确定度”、“下一最佳观测选哪个”（\cref{eq:gauss-mutual-info} 及 D-最优准则）。\cref{eq:ow-ib} 用的是\emph{同一个} $I(\cdot;\cdot)$，只是把它摆到一对\emph{新}的随机变量上：

\begin{itemize}
  \item 第一项 $I(\mathcal{X};\tilde{\mathcal{X}})$——\emph{压缩项}：压缩表示 $\tilde{\mathcal{X}}$ 还保留了多少关于原始基元 $\mathcal{X}$ 的信息。它越小，表示越精简（物体越少、越粗）。最小化它，就是在\emph{挤掉冗余}。
  \item 第二项 $I(\tilde{\mathcal{X}};\mathcal{Y})$——\emph{保真项}：压缩表示还保留了多少关于\emph{任务} $\mathcal{Y}$ 的信息。它越大，表示对完成任务越有用。\emph{最大化}它（故前面带负号、再乘 $\beta$），就是在\emph{留住任务要紧的区分}。
\end{itemize}

\noindent 于是 \cref{eq:ow-ib} 一句话即“\textbf{在所有把基元聚成物体的方案里，挑那个最精简、又最不损害任务的}”。$\beta$ 大则偏保真（物体切得细），$\beta$ 小则偏压缩（物体并得粗）——一架钢琴在“搬”任务下 $\mathcal{Y}$ 只关心整体、保真项不要求细分，于是压成一个物体；在“调音”任务下 $\mathcal{Y}$ 关心每根弦、保真项逼着细分，于是裂成上百个。\textbf{“该切多细”被 $\mathcal{Y}$（任务）与 $\beta$ 决定，而非拍脑袋。}这正是 \cref{tab:intro-eras} “任务驱动感知”的精确含义。
\end{insight}

\paragraph{怎么定义 $p(y\mid x)$、怎么解。}
信息瓶颈要可算，须先定义 $p(y\mid x)$——基元 $x$ 与任务 $y$ 的\emph{相关度}。Clio 用 CLIP 的图文余弦相似度作 $p(y\mid x)$ 的代理：把基元的图块特征与任务文本特征算 \cref{eq:ow-query} 式的相似度（这又回到 CLIP 共享空间这一支点）。求解上，\cref{eq:ow-ib} 可用\textbf{凝聚式信息瓶颈}（Agglomerative IB）\cite{slonim1999aib}\emph{增量}地解：建一张图，节点为基元、边连相邻基元，再\emph{逐步合并}。每条边 $(i,j)$ 赋一个权重 $d_{ij}$，度量合并基元簇 $\tilde{x}_i,\tilde{x}_j$ 所造成的\emph{信息失真}：
\begin{equation}
  d_{ij}=\big(p(\tilde{x}_i)+p(\tilde{x}_j)\big)\cdot D_{\mathrm{JS}}\!\big[p(y\mid\tilde{x}_i),\ p(y\mid\tilde{x}_j)\big],
  \label{eq:ow-aib}
\end{equation}
其中 $p(\tilde{x}_i)$ 直觉上记录簇 $i$ 已并入的基元数，$D_{\mathrm{JS}}$ 是\emph{Jensen--Shannon 散度}。每步贪心地合并权重最小的边——即合并“任务相关分布最像、合并后信息损失最小”的两簇。结果是把相邻基元逐步并成\emph{任务相关}的大物体。Clio 还把这一思路推到场景层级：把自由空间的“地点基元”聚成任务相关的区域（如房间），得到层为“物体基元--物体--地点基元--区域”的分层场景图。

\begin{insight}[$D_{\mathrm{JS}}$ 与本书的 KL 散度同源]
\cref{eq:ow-aib} 里的 Jensen--Shannon 散度不是新面孔：它就是本书 \cref{eq:gauss-kl} 那个 KL 散度的\emph{对称化}版本——$D_{\mathrm{JS}}(p\|q)=\tfrac12\mathrm{KL}(p\|m)+\tfrac12\mathrm{KL}(q\|m)$，$m=\tfrac12(p+q)$。KL 度量“两个分布有多不同”（\cref{eq:gauss-kl} 用于比较近似后验与真后验），$D_{\mathrm{JS}}$ 取其对称化、值域有界，正适合当“合并两簇的代价”。于是 \cref{eq:ow-aib} 一句话：\textbf{两簇的任务相关分布越像（$D_{\mathrm{JS}}$ 越小）、加权后越该先合并}——贪心地从最像的并起，正是凝聚式信息瓶颈对 \cref{eq:ow-ib} 的近似求解。本节从头到尾用的，是本书 \cref{sec:est-mahalanobis} 早已备好的同一套信息论工具（熵、互信息、KL/JS）。
\end{insight}

\begin{pitfall}[信息瓶颈的两处误差来源]
\cref{eq:ow-ib} 是一个\emph{一般}的任务驱动建图目标，但落到实处有两个误差源\cite{paull2026openworld}：(1) 用 CLIP 图文相似度\emph{近似} $p(y\mid x)$ 本身不准（CLIP 不懂否定与复杂语义）；(2) 跨视图\emph{聚合}语义知识的方式会引入误差。这两点都是当前活跃研究——例如有工作改用高斯泼溅地图、并对这些概率给出更扎实的估计。把 \cref{eq:ow-ib} 当成\emph{框架}（“任务驱动该被写成信息瓶颈”）是稳的；把某个具体实现里的 $p(y\mid x)$ 当成精确真值则不稳。
\end{pitfall}

% ════════════════════════════════════════════════════════════════
\section{展望：还需要显式地图吗？}
\label{sec:openworld-future}

\paragraph{动机：在大模型时代重问开篇之问。}
本书开篇（\cref{sec:intro-need}）追问过一个根本问题——\emph{自主机器人“真的”需要 SLAM 吗}？当时引 Cadena 等的回答，关键词是\emph{回环}：放弃回环，SLAM 就退化为里程计；SLAM 之所以有价值，在于它构建\emph{全局一致}的表示、抵御感知混叠、并满足许多任务对一致地图的硬需求。如今，基础模型从互联网级数据学到了惊人的视觉表征，于是有理由\emph{重新追问}同一问题：在这样的模型面前，还需不需要\emph{显式}地图？本节梳理这一开放问题的\emph{事实面}，最后给出本书自己的判断。

\begin{note}[读法提示：本节是综述，不是某派宣言]
本节话题（“还要不要地图”“VLA 会不会取代 SLAM”）在文献里立场鲜明、争论激烈。本书的处理是：\textbf{只陈述可复核的事实（基准、数据集、模型谱系），不替任何一派代言}。读到“某基准上 A 胜过 B”这类话，请理解为\emph{在那个特定基准上的实验结果}，而非“某某断言 A 永远更好”。本书自己的判断只在本节末以本书口吻给出，且明确标为本书立场。
\end{note}

\subsection{把地图喂给大模型（选读）}
\label{sec:openworld-mapapi}

一个自然方向是\emph{让 LLM 与地图表示更深地融合}\cite{paull2026openworld}。已有两条具体路径：

\emph{地图提示}（Map Prompting）：直接把地图以\emph{文本}形式塞进 LLM 的提示。当场景图的物体与边都用自然语言描述时（\cref{sec:openworld-objgraph}），整张图可写成结构化文本（如 JSON/YAML）加进提示。稀疏几何描述（物体质心、尺寸、框坐标）原则上也能加，但 LLM 能多好地利用这类几何信息尚不清楚。这种“VLM 造图、文本喂给 LLM”的模块化系统，可看作把若干预训练模型\emph{组合}起来用。

\emph{地图 API}（Map API）：文本提示传不了精细几何、且只适用于稀疏可文本化的地图。另一条路是把地图当成\emph{独立模块}、只通过\emph{符号化 API} 暴露功能——如 \texttt{exists(object)}、\texttt{where\_is(object)}、\texttt{distance(A,B)}：输入用自然语言、函数实现靠 \cref{sec:openworld-mapping} 的查询机制。LLM 调用可与地图函数调用、机器人 API 调用交错，搭出复杂系统（即让 LLM 在生成答案时“调用外部工具”）。

\subsection{重问：还需要显式地图吗?}
\label{sec:openworld-needmap}

围绕“显式地图能否被基础模型直接替代”，文献给出了若干\emph{实验事实}：

\begin{itemize}
  \item \textbf{长上下文 VLM 直喂图像帧}。新一代 VLM 能直接处理一长串 RGB 图像——本可拿去建图的那些帧。在具身问答基准 OpenEQA 上，\emph{直接}用 $50$ 帧提示 GPT-4V，其表现\emph{优于}在该基准上使用显式场景图提示\cite{paull2026openworld}。在导航场景下，Mobility VLA 用一整段办公室漫游（$948$ 帧）提示 VLM、要它根据用户查询找出目标图像、再交给基于地图的导航流水线；其作者还试过一个\emph{无地图}变体（让 VLM 直接根据当前观测与漫游输出导航动作），结果\emph{无效}\cite{paull2026openworld}。
  \item \textbf{几何/物理感知的 VLM}。多数 VLM 开箱缺乏对几何与物理属性的感知；有工作据此增强基础模型，使其能从\emph{单张}相机图像就对三维空间做相当不错的推理。但在 OpenEQA 这类基准上，纯感知增强的 VLM 与\emph{在流水线里纳入 SLAM} 所展现的具身智能之间，仍有明显差距\cite{paull2026openworld}。
\end{itemize}

\noindent 把这些事实并置，可读出一条经验规律：\emph{是否需要显式地图，似乎主要取决于任务的空间与时间跨度}——短跨度、单视角的问答，长上下文 VLM 可能够用；长跨度、需全局一致与回环纠错的任务，显式地图（SLAM）的价值依旧。这与开篇 \cref{sec:intro-need} 的结论\emph{并不矛盾、而是延续}：SLAM 的不可替代性，恰恰在“全局一致、抗感知混叠、长时一致”这些\emph{长跨度}诉求上——大模型擅长的是另一端。

\subsection{机器人基础模型与 VLA（选读）}
\label{sec:openworld-vla}

语言与二维视觉的基础模型之所以好建，在于有互联网级数据可训。一个自然的追问是：能否为\emph{机器人}也造一个基础模型——一个能跨多种机器人本体、执行多样技能的\emph{通用策略}？这把问题引向本书前言承诺的另一条线：\emph{视觉–语言–动作}（VLA）模型与具身智能。这里梳理其\emph{事实面}。

\paragraph{数据集。}
机器人基础模型的瓶颈是数据——真机采集极其昂贵。代表性努力\cite{paull2026openworld}：RT-1 的数据集含 $13$ 台机器人在 $17$ 个月里执行 $700$ 种任务的 $13$ 万条轨迹；BridgeData V2 增其任务多样性；DROID 是多机构协作的操作数据集；这些与更多数据被整合成 \textbf{Open-X-Embodiment}——含 $22$ 种机器人本体、$527$ 种技能、超百万条片段，跨 $34$ 个实验室。高效的人类示教接口（VR 遥操作、低成本双臂 ALOHA、与机器人夹爪精确匹配的手持采集器 UMI）是规模化采集的关键。这些数据大多形如\emph{传感器-动作}轨迹 $\mathcal{D}=\{(\mathbf{z}_t,\mathbf{u}_t)\}_{t=1}^{T}$——主要瞄准学习\emph{技能}（抓取、操作），复杂任务则靠串联子技能或语言条件化来达成。

\paragraph{模型谱系。}
最朴素的学法是\emph{行为克隆}（behavior cloning）：监督地学“观测 $\to$ 动作”的映射——但已知它脆弱、易发散（示教里少有次优数据，动作稍偏离示教就出分布）。故大规模数据上训出的“机器人基础模型”候选，多沿用 \cref{sec:openworld-fm} 介绍的基础模型结构，只是把 Transformer 的输出 token 直接解读为\emph{控制动作}——这便是 \textbf{VLA}（Vision-Language-Action）模型\cite{paull2026openworld}。谱系上的代表：GATO（单一模型靠统一 tokenize 跨任务，含机器人控制）、PaLM-E（具身多模态语言模型，混合视觉-语言与操作数据训练）、RT-1/RT-2（及在 Open-X-Embodiment 上训的 RT-1-X/RT-2-X）、ACT（动作分块 Transformer，一次输出\emph{一段}动作序列而非单步，对精细操作尤要）、OpenVLA（全开源组件搭建）、$\pi_0$。另一条不同范式是\emph{扩散策略}（diffusion policy）：把视觉-运动策略表示成一个\emph{条件去噪}过程，从随机噪声迭代细化出动作。两条范式并非互斥——Octo 就用大 Transformer 产出、再喂给一个跑前向扩散的“动作头”。\emph{强化学习}（RL）也日渐流行，常用作在仿真中微调上述策略（仿真里奖励好定义、试错成本低）。

\begin{note}[把 VLA 接回本书：它学的是策略，不是地图]
VLA 与本书主体（状态估计、SLAM）处理的是\emph{不同}的问题：本书前面各章解决“\emph{感知与表示}世界”（我在哪、世界长什么样），VLA 解决“\emph{在世界里如何行动}”（给定观测该出什么动作）。前言把全书弧线写作“经典模型 $\to$ 学习方法 $\to$ 具身智能”，VLA 正是“具身智能”这一端。它们在数据形式上甚至相通——VLA 的轨迹 $\{(\mathbf{z}_t,\mathbf{u}_t)\}$ 里，$\mathbf{u}_t$ 是\emph{动作}（控制量），这与本书运动模型里的控制输入 $\mathbf{u}$ 是同一类量。本书不展开 VLA 的训练细节（那属另一部专著），此处只确立其在全书版图里的\emph{位置}：地图回答“世界是什么”，策略回答“该怎么动”，二者是具身智能的两半。
\end{note}

\subsection{收束：本书的判断}
\label{sec:openworld-closing}

把本节的事实面收拢，全书“是否需要 SLAM / 地图”这条贯穿线（\cref{sec:intro-need} 提出、本章重问）可以收尾了。文献里有两种可能的图景\cite{paull2026openworld}：一种是\emph{进一步扩大}机器人数据与专用基础模型、结合互联网级 VLM/LLM，或许能使显式地图在\emph{某些}场景下\emph{不再必要}；另一种（在大尺度、组合性任务上同样可能、甚至更可能）是\emph{真正的泛化与可扩展}仍需某种\emph{显式结构}——而这种结构恰可由 SLAM 这样的过程学到。

\begin{insight}[本书立场：互补，而非取代——与开篇首尾相扣]
本书在此明确给出自己的判断（区别于上面陈述的文献事实）：\textbf{显式世界建模（SLAM）与通用机器人策略（VLA），最可能是\emph{互补}的两半，而非彼此取代。}理由正是本书一路铺陈下来的：
\begin{itemize}
  \item \emph{感知、表示、理解世界，是采取行动的前提}（至少在本书关心的复杂问题里）。一个能可靠地建出全局一致地图、抵御感知混叠、做长时一致推理的系统，只会\emph{帮助}而非妨碍鲁棒地行动。这与开篇 \cref{sec:intro-need} “SLAM 是抵御错误数据关联的天然防线”一脉相承。
  \item 本章的事实也支持这一判断：长跨度、需全局一致的任务里，纳入 SLAM 的流水线仍显著领先纯 VLM（\cref{sec:openworld-needmap}）；连 map-free 的导航变体也被验证\emph{无效}。是否需要显式地图\emph{取决于任务的空间与时间跨度}——而非“地图已被取代”。
\end{itemize}
\noindent 于是全书的弧线在此闭合：从 \cref{sec:intro-need} 问“要不要 SLAM”，经整部经典 SLAM 把“怎么建一致的几何地图”讲透，经这一部把“怎么让地图学习、带语义、对开放世界开放”讲透，最终在大模型时代\emph{重问}并\emph{重答}同一问题——答案不是“地图过时了”，而是“地图与策略，各司其职、合则两利”。这正是前言所谓“架起一座由传统机器人学通向现代具身智能的桥梁”：桥的此岸是几何与一致性（本书主体），彼岸是通用策略与行动（VLA），而本章是桥面合龙的最后一段。更完整的收尾寄语，见全书结语（\cref{ch:epilogue}）。
\end{insight}

% ════════════════════════════════════════════════════════════════
\section*{常见误解}

\begin{pitfall}[误解一：“开放词汇前端”就等于“开放世界地图”]
不等。把前端检测器换成能接受任意文本的开放词汇版本，得到的地图\emph{仍是闭集}——因为“要找哪些概念”在\emph{建图时}就被那次文本提示锁死了（\cref{sec:openworld-closed-open}）。真正的开放世界地图，是把开放词汇\emph{特征} $\boldsymbol\psi_k$ 嵌进地图、把“是不是某概念”的判断\emph{推迟到用图时}（\cref{eq:ow-query}），且查询集 $\mathsf{C}$ 可随时重定义。判别准绳只有一句：\emph{概念集定死在建图时，还是推迟到查询时？}
\end{pitfall}

\begin{pitfall}[误解二：开放词汇 SLAM 是一套全新的、不依赖经典 SLAM 的系统]
恰相反。本章绝大多数方法严格说是\emph{纯建图}：它们假设已有上游 SLAM（前面各章）提供\emph{带位姿}的 RGB-D 帧，自己只负责往几何骨架上\emph{嵌一层特征}（\cref{eq:ow-mapelem} 的几何-语义解耦）。几何 $\mathbf{p}_k$ 仍是经典 SLAM 解出来的；本章接管的只是\emph{语义分量} $\boldsymbol\psi_k$。把它当成“取代经典 SLAM 的新东西”，就既看不清它站在谁的肩上，也理解不了它为何能不重做几何。
\end{pitfall}

\begin{pitfall}[误解三：信息瓶颈 \cref{eq:ow-ib} 是个深奥的新理论]
不是。它用的就是本书 \cref{sec:est-mahalanobis} 早已定义的\emph{互信息} $I(\cdot;\cdot)$（\cref{eq:mutual-info}）——那里用它做主动感知、选下一最佳观测。\cref{eq:ow-ib} 只是把同一个 $I$ 摆到“基元 $\mathcal{X}$ / 压缩表示 $\tilde{\mathcal{X}}$ / 任务 $\mathcal{Y}$”这对新变量上：压缩项 $I(\mathcal{X};\tilde{\mathcal{X}})$ 要小（精简）、保真项 $I(\tilde{\mathcal{X}};\mathcal{Y})$ 要大（保任务）。求解用的 $D_{\mathrm{JS}}$（\cref{eq:ow-aib}）也只是本书 \cref{eq:gauss-kl} 那个 KL 散度的对称化。把它读成本书信息论语言的又一次落地，远比当成天书有用。
\end{pitfall}

\begin{pitfall}[误解四：物体级地图“一物一向量”就是单纯的有损压缩，没别的好处]
不止省内存。一物一向量的物体表示，天然契合具身 AI 的\emph{物体检索}（“找点能喝的”$\to$ 直接返回汽水的点云交给规划）与\emph{以物体为中心的任务规划}，还能往上聚成区域/楼层做\emph{分层语言导航}（\cref{eq:ow-hsg}）。粒度自低（稠密）向高（场景图）是\emph{用表达力换内存与可检索性}——而“到底该换到哪一档”，正是 \cref{sec:openworld-task} 信息瓶颈要回答的任务驱动问题。
\end{pitfall}

\begin{pitfall}[误解五：大模型能直喂图像帧，所以 SLAM/显式地图过时了]
这是把\emph{特定基准上的实验结果}误读成普遍结论。事实是：短跨度问答里长上下文 VLM 可能够用，但长跨度、需全局一致的任务里，纳入 SLAM 的流水线仍显著领先，连 map-free 导航变体都被验证无效（\cref{sec:openworld-needmap}）。\emph{是否需要显式地图取决于任务的空间与时间跨度}。本书的判断（\cref{sec:openworld-closing}）是显式地图与通用策略\emph{互补}——这与开篇 \cref{sec:intro-need} 对 SLAM 价值的论断首尾一致，而非被它推翻。
\end{pitfall}

\begin{pitfall}[误解六：把语义特征 $\boldsymbol\psi$ 与视线方向 $\mathbf{f}$、光流 $\mathbf{f}^{\mathrm{flow}}$ 混为一谈]
本章特意\emph{不}用 $\mathbf{f}$ 记特征——因 \cref{ch:dense} 已用 $\mathbf{f}$ 表“单位视线方向”、\cref{ch:learning} 已用 $\mathbf{f}^{\mathrm{flow}}$ 表光流场。本章一切语义特征/嵌入一律记 $\boldsymbol\psi$（已在“本章记号与约定”框登记），编码器记 $\mathcal{E}_X$。读公式时认准 $\boldsymbol\psi$，免得把“语义特征”当成“几何视线”。
\end{pitfall}

% ════════════════════════════════════════════════════════════════
\section*{本章小结}

本章是全书“经典 $\to$ 学习 $\to$ 开放世界”弧线（\cref{sec:intro-frontier}、前言）在内容上的落点，也是“学习时代”这一部的收束。一条主线贯穿始终：\textbf{从闭集语义（\cref{ch:semantic}）到开放世界，要换掉的只是地图元素里“是什么”那一步——把有限词典里的标签 $\sigma$ 换成基础模型给出的连续特征 $\boldsymbol\psi$（\cref{eq:ow-mapelem}），并把“是不是某概念”的判断从建图时推迟到查询时}。这一替换之上，开放词汇建图沿“$\mathrm{Seg}\to$ 抽特征 $\to$ 反投影 $\to$ 关联融合（\cref{eq:ow-fusion,eq:ow-objfuse}）$\to$ 余弦查询（\cref{eq:ow-query}）”一条内核展开，在稠密、物体级/场景图、隐式三类几何上变体——分歧只在特征存到哪一级几何、是\emph{用表达力换内存与可检索性}。

本章的数学重点是\emph{任务驱动的地图压缩}：“场景该切成几个物体”这一开集时代独有的病态问题，被\emph{信息瓶颈}（\cref{eq:ow-ib}）形式化为一个由任务 $\mathcal{Y}$ 与权衡 $\beta$ 决定的优化——压缩项要小、保任务信息项要大。它用的正是本书 \cref{sec:est-mahalanobis} 早备好的互信息 $I(\cdot;\cdot)$（\cref{eq:mutual-info}）与 KL/JS 散度（\cref{eq:gauss-kl}），是“鲁棒感知时代”（\cref{tab:intro-eras}）\emph{任务驱动感知}诉求的精确落地。

最后，本章在大模型时代\emph{重问}开篇 \cref{sec:intro-need} 的“是否需要 SLAM”：梳理长上下文 VLM 与 VLA 的事实后，给出本书自己的判断——显式世界建模（SLAM）与通用机器人策略（VLA）是\emph{互补}的两半，是否需要显式地图\emph{取决于任务的空间与时间跨度}。全书弧线在此闭合：地图回答“世界是什么”、策略回答“该怎么动”，合则两利。更完整的收尾见结语（\cref{ch:epilogue}）。

\begin{table}[H]\centering\small
\caption{符号表}
\begin{tabular}{lll}
\toprule
符号 & 含义 & 首次出现 \\
\midrule
$\boldsymbol\psi\in\mathbb{R}^{D}$ & 语义特征/嵌入（避让视线 $\mathbf{f}$、光流 $\mathbf{f}^{\mathrm{flow}}$）& \cref{sec:openworld-feature} \\
$\mathcal{E}_X,\ \mathcal{E}_{\mathrm{img}},\ \mathcal{E}_{\mathrm{txt}}$ & 模态 $X$/图像/文本编码器 & \cref{eq:ow-encoder} \\
$\langle\boldsymbol\psi_a,\boldsymbol\psi_b\rangle$ & 余弦相似度（$\boldsymbol\psi$ 已单位化）& \cref{sec:openworld-feature} \\
$m_k=(\mathbf{p}_k,\boldsymbol\psi_k)$ & 地图元素：几何 $+$ 特征；$\mathcal{M}=\{m_k\}$ & \cref{eq:ow-mapelem} \\
$\mathrm{Seg}(I)=\{Z_i\}$ & 类无关分割掩膜（SAM）& \cref{sec:openworld-sam} \\
$T,\ \boldsymbol\psi_T,\ \mathsf{C}$ & 文本查询/其特征/候选类别集 & \cref{eq:ow-clip-cls} \\
$\boldsymbol\psi^{G},\boldsymbol\psi_i^{L},\boldsymbol\psi_{u,v}^{P}$ & 全局/局部/逐像素混合特征 & \cref{sec:openworld-dense} \\
$c_k,\ \alpha$ & 元素累积置信/单次观测置信 & \cref{eq:ow-fusion} \\
$o_j=(\mathsf{P}_{o_j},\boldsymbol\psi_{o_j},n_{o_j})$ & 物体：点云/特征/检测计数；$\mathcal{V}_O$ & \cref{sec:openworld-objgraph} \\
$\phi_{\mathrm{geo}},\phi_{\mathrm{sim}}$ & 几何/语义相似度（$\in[0,1]$）& \cref{eq:ow-objfuse} \\
$\mathcal{G}=(\mathcal{V},\mathcal{E})$ & （分层）场景图 & \cref{eq:ow-hsg} \\
$\mathcal{X},\tilde{\mathcal{X}},\mathcal{Y}$ & 基元集/压缩表示/任务集 & \cref{eq:ow-ib} \\
$I(\cdot;\cdot),\ \beta$ & 互信息（同 \cref{eq:mutual-info}）/权衡系数 & \cref{eq:ow-ib} \\
$D_{\mathrm{JS}}$ & Jensen--Shannon 散度（$\mathrm{KL}$ 对称化，\cref{eq:gauss-kl}）& \cref{eq:ow-aib} \\
$\mathcal{D}=\{(\mathbf{z}_t,\mathbf{u}_t)\}$ & VLA 的传感器-动作轨迹 & \cref{sec:openworld-vla} \\
\bottomrule
\end{tabular}
\end{table}

\begin{table}[H]\centering\small
\caption{定理 / 公式速查表}
\begin{tabular}{lll}
\toprule
公式 & 一句话说明 & 对应 \\
\midrule
地图元素 \cref{eq:ow-mapelem} & 闭集标签 $\sigma$ 换成连续特征 $\boldsymbol\psi$ & \cref{sec:openworld-closed-open} \\
CLIP 零样本分类 \cref{eq:ow-clip-cls} & 图文同空间比余弦；后文查询母型 & \cref{sec:openworld-clip} \\
ConceptFusion 融合 \cref{eq:ow-fusion} & 置信加权递归平均；同 TSDF/卡尔曼智慧 & \cref{sec:openworld-dense} \\
零样本查询 \cref{eq:ow-query} & 地图特征对文本特征取余弦近邻；$\mathsf{C}$ 可重定义 & \cref{sec:openworld-dense} \\
物体融合 \cref{eq:ow-objfuse} & 计数加权融特征、并点云 & \cref{sec:openworld-objgraph} \\
分层场景图 \cref{eq:ow-hsg} & 楼层-区域-物体分层、逐层语言定位 & \cref{sec:openworld-objgraph} \\
LERF 相关性分数 \cref{eq:ow-lerf} & 余弦查询 $+$ 规范短语去噪（渲染嵌入上）& \cref{sec:openworld-implicit} \\
信息瓶颈 \cref{eq:ow-ib} & 任务驱动该切多细：压缩 vs 保任务 & \cref{sec:openworld-task} \\
凝聚式 IB 边权 \cref{eq:ow-aib} & 任务分布越像（$D_{\mathrm{JS}}$ 小）越先合并 & \cref{sec:openworld-task} \\
\bottomrule
\end{tabular}
\end{table}

\begin{table}[H]\centering\small
\caption{知识点总表}
\begin{tabular}{clll}
\toprule
\# & 知识点 & 核心要点 & 难度 \\
\midrule
1 & 闭集 $\to$ 开集 & 标签换特征，判断推迟到查询时 & $\bigstar\bigstar$ \\
2 & 两种“开放”之辨 & 开放词汇前端仍闭集；开放词汇地图才真开放 & $\bigstar\bigstar$ \\
3 & CLIP 共享空间 & 图文同尺比余弦，查询支点 & $\bigstar\bigstar$ \\
4 & 类无关分割 SAM & 抠区域以抽细粒度特征 & $\bigstar$ \\
5 & range-feature 观测 & 测距点带特征；range-category 的连续化 & $\bigstar\bigstar$ \\
6 & 特征融合与查询 & 置信加权平均；余弦近邻；$\mathsf{C}$ 可重定义 & $\bigstar\bigstar\bigstar$ \\
7 & 物体级与场景图 & 一物一向量；分层语言定位 & $\bigstar\bigstar\bigstar$ \\
8 & 隐式开放词汇（选读）& LERF/LangSplat；同一内核连续化 & $\bigstar\bigstar\bigstar$ \\
9 & 任务驱动 = 信息瓶颈 & 压缩 vs 保任务；接本书互信息 & $\bigstar\bigstar\bigstar\bigstar$ \\
10 & 凝聚式 IB 求解 & $D_{\mathrm{JS}}$ 贪心合并；接本书 KL & $\bigstar\bigstar\bigstar\bigstar$ \\
11 & 是否还需要地图 & 取决任务跨度；与 \cref{sec:intro-need} 呼应 & $\bigstar\bigstar$ \\
12 & VLA 与具身智能（选读）& 策略学行动，与地图互补 & $\bigstar\bigstar$ \\
\bottomrule
\end{tabular}
\end{table}

% ════════════════════════════════════════════════════════════════
\section*{延伸阅读}
\begin{itemize}
  \item \textbf{本章蓝本}：SLAM Handbook\cite{paull2026openworld} 第 17 章《Towards Open-World Spatial AI》——基础模型、开放词汇建图、任务驱动表示与“是否需要地图”的系统综述（本章框架与系统选择的出处）。
  \item \textbf{基础模型}：CLIP\cite{radford2021clip}（图文对比对齐，本章查询支点）、SAM\cite{kirillov2023sam}（类无关分割）、DINO/DINOv2\cite{caron2021dino,oquab2024dinov2}（自监督视觉特征）；深度学习通论\cite{goodfellow2016deeplearning}。
  \item \textbf{开放词汇建图}：稠密 ConceptFusion\cite{jatavallabhula2023conceptfusion}；物体级/场景图 ConceptGraphs\cite{gu2024conceptgraphs}、分层 HOV-SG\cite{werby2024hovsg}；隐式 LERF\cite{kerr2023lerf}、LangSplat\cite{qin2024langsplat}。
  \item \textbf{任务驱动表示}：Clio\cite{maggio2024clio}（信息瓶颈式任务驱动建图）；信息瓶颈原理\cite{tishby1999ib}、凝聚式信息瓶颈\cite{slonim1999aib}。本书的互信息/KL 见 \cref{sec:est-mahalanobis}。
  \item \textbf{VLA 与机器人基础模型}：综述见本章蓝本\cite{paull2026openworld} §17.4；具体模型与数据集（RT/PaLM-E/OpenVLA/Open-X-Embodiment 等）原文从略，属另一专题。
\end{itemize}

\section*{本章与前后章节的关系}
\begin{table}[H]\centering\small
\begin{tabularx}{\linewidth}{lLL}
\toprule
章节 & 关系 & 衔接点 \\
\midrule
\cref{ch:intro} 初识 SLAM & 提供弧线与“是否需要 SLAM”之问 & \cref{sec:intro-need} 与本章末首尾相扣；\cref{tab:intro-eras} 任务驱动诉求 \\
\cref{ch:slam_est} 状态估计 & 提供互信息/熵/KL 的本书定义 & \cref{eq:mutual-info,eq:gauss-kl} 是信息瓶颈 \cref{eq:ow-ib} 的前置 \\
\cref{ch:dense} 稠密建图 & 提供几何表示与逐元素融合 & 特征嵌进 \cref{ch:dense} 的稠密表示；\cref{eq:dense-octo-logodds} 是 range-feature 之祖 \\
\cref{ch:learning} 学习赋能 & 立“学习只换母方程零件”之统一立场 & 本章把“被接管的零件”推到语义层；ViT/CLIP 主干同源 \\
\cref{ch:semantic} 度量-语义 & 闭集语义、range-category、场景图 & 本章直接接住其闭集边界，标签 $\sigma\to$ 特征 $\boldsymbol\psi$ \\
\cref{ch:dynamic} 动态可变形 & 另一条“非经典假设”前沿线 & 本章假设静态场景；动态语义留给 \cref{ch:dynamic} \\
\cref{ch:epilogue} 结语 & 全书收尾寄语 & 本章末“互补”判断在结语合流 \\
\bottomrule
\end{tabularx}
\end{table}

% ════════════════════════════════════════════════════════════════
\section*{习题}
\begin{enumerate}
  \item \textbf{闭集到开集的本质。}\;用一句话说清：从 \cref{ch:semantic} 的闭集语义地图到本章的开放词汇地图，地图元素 $m_k$ 里\emph{哪个}分量被替换、换成了什么、以及“是不是某概念”的判断从\emph{何时}挪到了\emph{何时}。再举一个闭集词典必然失灵、而开放词汇能处理的查询例子。
  \item \textbf{两种“开放”之辨。}\;有人把一台能接受任意文本提示的开放词汇检测器接到 SLAM 前端，宣称“这就是开放世界 SLAM”。用本章的判别准绳反驳他：这样得到的地图为什么\emph{骨子里仍是闭集}？要让它真正开放，该把什么\emph{嵌进}地图、把什么判断\emph{推迟}到何时？
  \item \textbf{查询即 CLIP 分类。}\;写出地图查询 \cref{eq:ow-query} 与 CLIP 图像分类器 \cref{eq:ow-clip-cls}，逐项对应：哪个量从“整图特征”换成了“地图元素特征”？并解释为什么“候选类别集 $\mathsf{C}$ 不必在建图时已知”正是\emph{开集}相对\emph{闭集}的根本自由。
  \item \textbf{融合即信息加权。}\;说明 ConceptFusion 的特征融合 \cref{eq:ow-fusion} 为何是一个“权重递增的加权平均”，并把它与 \cref{ch:dense} 的 TSDF 加权融合、卡尔曼更新“按信息量加权”对应起来（提示：把 $c_k$ 读作已累积信息、$\alpha$ 读作本次信息）。再说明为什么有方法宁可取\emph{中位元}也不平均高维特征。
  \item \textbf{粒度的病态与信息瓶颈。}\;用“搬钢琴 / 弹钢琴 / 调音”三例说明“场景该切几个物体”为何在\emph{指定任务前}无解。再对照信息瓶颈 \cref{eq:ow-ib}：当任务 $\mathcal{Y}$ 从“搬”变成“调音”时，\emph{保真项} $I(\tilde{\mathcal{X}};\mathcal{Y})$ 如何逼着物体粒度变细？$\beta$ 增大又会把粒度推向哪一端？
  \item \textbf{信息瓶颈接本书互信息。}\;指出 \cref{eq:ow-ib} 里的 $I(\cdot;\cdot)$ 与本书 \cref{eq:mutual-info} 是同一个量。用一句话各解释\emph{压缩项} $I(\mathcal{X};\tilde{\mathcal{X}})$ 与\emph{保真项} $I(\tilde{\mathcal{X}};\mathcal{Y})$ 的含义；再说明 \cref{eq:ow-aib} 的 $D_{\mathrm{JS}}$ 与本书 \cref{eq:gauss-kl} 的 KL 散度是什么关系，以及“为什么 $D_{\mathrm{JS}}$ 小的两簇该先合并”。
  \item \textbf{三类表示的取舍。}\;为下列三个任务各挑\emph{稠密 / 物体级 / 分层场景图}中最合适的一种表示并说明理由：(a) 移动机器人要“去二楼会议室拿那支红笔”；(b) 桌面机械臂要分辨并抓取一堆杂物里某个特定零件；(c) AR 应用要对整个房间做最细粒度的语义渲染。把你的选择对应到“表达力 vs 内存/可检索性”的权衡上。
  \item \textbf{是否还需要地图。}\;复述 \cref{sec:openworld-needmap} 给出的两条实验事实（OpenEQA 直喂 $50$ 帧、Mobility VLA 的 map-free 变体），并据此论证本书的判断“是否需要显式地图取决于任务的空间与时间跨度”。再把这一判断接回开篇 \cref{sec:intro-need}：它\emph{推翻}还是\emph{延续}了“放弃回环 SLAM 就退化为里程计”的论断？为什么？
  \item \textbf{（选读）地图 API 与工具调用。}\;设计三个你认为最有用的符号化地图 API 函数（除正文给的 \texttt{exists/\allowbreak where\_is/\allowbreak distance} 外），各写出输入输出，并说明 LLM 在生成一个机器人计划时会\emph{如何交错}调用它们与机器人动作 API。再说明“地图 API”相比“把地图写成文本塞进提示”，在传递\emph{几何}信息上的优势。
  \item \textbf{（选读）VLA 在全书版图里的位置。}\;用本章 \cref{sec:openworld-vla} 的“注记”说清：VLA 解决的问题与本书前面各章（状态估计、SLAM）解决的问题\emph{有何不同}？VLA 轨迹 $\{(\mathbf{z}_t,\mathbf{u}_t)\}$ 里的 $\mathbf{u}_t$ 与本书运动模型的控制输入 $\mathbf{u}$ 是否同一类量？据此论证“地图回答世界是什么、策略回答该怎么动，二者互补”这一全书收尾判断。
\end{enumerate}
